% paper68-technical-debt.tex — Measuring and Managing Technical Debt
%   via Trust Degradation and Maturity Assessment
%   Paper 68 of the Judgment Geometry series.
% Compile: pdflatex paper68-technical-debt
\documentclass[11pt]{article}
\usepackage{lmodern}
% jugeo-common.tex — shared preamble for all Judgment Geometry papers
% Include via: % jugeo-common.tex — shared preamble for all Judgment Geometry papers
% Include via: % jugeo-common.tex — shared preamble for all Judgment Geometry papers
% Include via: \input{jugeo-common}

% ── Packages ────────────────────────────────────────────────────────────
\usepackage{amsmath,amssymb,amsthm,mathtools}
\usepackage{tikz-cd}
\usepackage{listings}
\usepackage{xcolor}
\usepackage{xspace}
\usepackage{booktabs}
\usepackage{multirow}
\usepackage{enumitem}
\usepackage[expansion=false]{microtype}
\usepackage{hyperref}
\usepackage{cleveref}
% \usepackage{thmtools}  % disabled: not available in this TeX installation
\usepackage{float}
\usepackage{caption}
\usepackage{subcaption}
\usepackage{tabularx}
\usepackage{stmaryrd}       % \llbracket, \rrbracket
\usepackage{bbm}            % \mathbbm{1}

% ── Page geometry ───────────────────────────────────────────────────────
\usepackage[margin=1in]{geometry}

% ── Theorem environments ────────────────────────────────────────────────
\theoremstyle{plain}
\newtheorem{theorem}{Theorem}[section]
\newtheorem{lemma}[theorem]{Lemma}
\newtheorem{proposition}[theorem]{Proposition}
\newtheorem{corollary}[theorem]{Corollary}

\theoremstyle{definition}
\newtheorem{definition}[theorem]{Definition}
\newtheorem{example}[theorem]{Example}
\newtheorem{construction}[theorem]{Construction}

\theoremstyle{remark}
\newtheorem{remark}[theorem]{Remark}
\newtheorem{notation}[theorem]{Notation}

% ── Python listings style ──────────────────────────────────────────────
\definecolor{jugeoBlue}{HTML}{2563EB}
\definecolor{jugeoGreen}{HTML}{059669}
\definecolor{jugeoRed}{HTML}{DC2626}
\definecolor{jugeoPurple}{HTML}{7C3AED}
\definecolor{jugeoGray}{HTML}{6B7280}
\definecolor{jugeoBg}{HTML}{F8FAFC}

\lstdefinestyle{jugeo-python}{
  language=Python,
  basicstyle=\ttfamily\small,
  keywordstyle=\color{jugeoBlue}\bfseries,
  stringstyle=\color{jugeoGreen},
  commentstyle=\color{jugeoGray}\itshape,
  numberstyle=\tiny\color{jugeoGray},
  backgroundcolor=\color{jugeoBg},
  numbers=left,
  numbersep=6pt,
  frame=single,
  framerule=0.4pt,
  rulecolor=\color{jugeoGray!40},
  breaklines=true,
  breakatwhitespace=true,
  showstringspaces=false,
  tabsize=4,
  xleftmargin=1.5em,
  framexleftmargin=1.5em,
  morekeywords={dataclass,frozen,field,Enum,IntEnum,auto,Any,
                Mapping,Sequence,Optional,match,case},
  literate={→}{{$\to$}}1 {←}{{$\leftarrow$}}1
           {≤}{{$\leq$}}1 {≥}{{$\geq$}}1 {≠}{{$\neq$}}1
           {∈}{{$\in$}}1 {∀}{{$\forall$}}1 {∃}{{$\exists$}}1
           {⊕}{{$\oplus$}}1 {⊖}{{$\ominus$}}1,
  escapeinside={(*@}{@*)},
}
\lstset{style=jugeo-python}

\lstdefinestyle{jugeo-output}{
  basicstyle=\ttfamily\small,
  backgroundcolor=\color{jugeoBg},
  frame=single,
  framerule=0.4pt,
  rulecolor=\color{jugeoGray!40},
  breaklines=true,
  showstringspaces=false,
  xleftmargin=1.5em,
  framexleftmargin=0.5em,
  numbers=none,
}

% ── JuGeo-specific macros ──────────────────────────────────────────────

% Judgment 8-tuple
\newcommand{\J}{\mathcal{J}}
\newcommand{\Jud}[8]{(#1,\,#2,\,#3,\,#4,\,#5,\,#6,\,#7,\,#8)}
\newcommand{\JudShort}[1]{\J_{#1}}

% Components
\newcommand{\coord}{\mathsf{c}}            % coordinate
\newcommand{\prop}{\varphi}                 % proposition
\newcommand{\carrier}{\mathcal{A}}          % carrier type
\newcommand{\evid}{\mathcal{E}}             % evidence bundle
\newcommand{\oblig}{\mathcal{O}}            % obligations
\newcommand{\obstr}{\mathcal{B}}            % obstructions
\newcommand{\trust}{\mathcal{T}}            % trust annotation
\newcommand{\prov}{\Pi}                     % provenance

% Trust algebra
\newcommand{\Talg}{\mathfrak{T}}            % trust ordered algebra
\newcommand{\Eadm}{\mathcal{E}_{\mathrm{adm}}}
\newcommand{\tleq}{\preceq}                 % trust ordering
\newcommand{\tjoin}{\oplus}                 % conservative join
\newcommand{\tatten}{\ominus}               % attenuation
\newcommand{\tpromote}{\uparrow_\pi}        % promotion
\newcommand{\tdemote}{\downarrow_\chi}       % demotion

% Trust tiers
\newcommand{\Tcontra}{\ensuremath{\mathsf{CONTRADICTED}}}
\newcommand{\Tunver}{\ensuremath{\mathsf{UNVERIFIED}}}
\newcommand{\Tcopilot}{\ensuremath{\mathsf{COPILOT}}}
\newcommand{\Toracle}{\ensuremath{\mathsf{ORACLE}}}
\newcommand{\Truntime}{\ensuremath{\mathsf{RUNTIME}}}
\newcommand{\Tsolver}{\ensuremath{\mathsf{SOLVER}}}
\newcommand{\Tproof}{\ensuremath{\mathsf{PROOF}}}

% Site and geometry
\newcommand{\Site}{\mathcal{S}}             % semantic site
\newcommand{\Cov}{\mathrm{Cov}}            % covering family
\newcommand{\Ob}{\mathrm{Ob}}              % objects
\newcommand{\Mor}{\mathrm{Mor}}            % morphisms
\newcommand{\res}{\mathrm{res}}            % restriction
\newcommand{\Cech}{\check{C}}              % Čech complex
\newcommand{\Hcoh}[1]{H^{#1}}             % cohomology group

% Descent
\newcommand{\descent}{\mathrm{desc}}
\newcommand{\glue}{\mathrm{glue}}
\newcommand{\overlap}[2]{#1 \cap #2}

% Semantic moves
\newcommand{\VERIFY}{\textsc{Verify}}
\newcommand{\CONSTRUCT}{\textsc{Construct}}
\newcommand{\NORMALIZE}{\textsc{Normalize}}
\newcommand{\GLUE}{\textsc{Glue}}
\newcommand{\RESTRICT}{\textsc{Restrict}}
\newcommand{\TRANSPORT}{\textsc{Transport}}
\newcommand{\REPAIR}{\textsc{Repair}}
\newcommand{\PROMOTE}{\textsc{Promote}}
\newcommand{\CHALLENGE}{\textsc{Challenge}}
\newcommand{\ESCALATE}{\textsc{Escalate}}

% Fragments
\newcommand{\QFLIA}{\texttt{QF\_LIA}}
\newcommand{\QFLRA}{\texttt{QF\_LRA}}
\newcommand{\QFBV}{\texttt{QF\_BV}}
\newcommand{\QFUF}{\texttt{QF\_UF}}

% Misc
\newcommand{\Python}{\textsc{Python}}
\newcommand{\JuGeo}{\textsc{JuGeo}}
\newcommand{\LEAN}{\textsc{Lean}\,4}
\newcommand{\Fstar}{F$^{\star}$}
\newcommand{\Coq}{\textsc{Coq}}
\newcommand{\Dafny}{\textsc{Dafny}}

% Common shortcuts
\newcommand{\ie}{i.e.\@\xspace}
\newcommand{\eg}{e.g.\@\xspace}
\newcommand{\cf}{cf.\@\xspace}
\newcommand{\wrt}{w.r.t.\@\xspace}

% Evaluation shortcuts
\newcommand{\numcases}{300}
\newcommand{\numspec}{100}
\newcommand{\numeq}{100}
\newcommand{\numbug}{100}
\newcommand{\medtime}{6.5\,ms}
\newcommand{\totaltime}{1.949\,s}


% ── Packages ────────────────────────────────────────────────────────────
\usepackage{amsmath,amssymb,amsthm,mathtools}
\usepackage{tikz-cd}
\usepackage{listings}
\usepackage{xcolor}
\usepackage{xspace}
\usepackage{booktabs}
\usepackage{multirow}
\usepackage{enumitem}
\usepackage[expansion=false]{microtype}
\usepackage{hyperref}
\usepackage{cleveref}
% \usepackage{thmtools}  % disabled: not available in this TeX installation
\usepackage{float}
\usepackage{caption}
\usepackage{subcaption}
\usepackage{tabularx}
\usepackage{stmaryrd}       % \llbracket, \rrbracket
\usepackage{bbm}            % \mathbbm{1}

% ── Page geometry ───────────────────────────────────────────────────────
\usepackage[margin=1in]{geometry}

% ── Theorem environments ────────────────────────────────────────────────
\theoremstyle{plain}
\newtheorem{theorem}{Theorem}[section]
\newtheorem{lemma}[theorem]{Lemma}
\newtheorem{proposition}[theorem]{Proposition}
\newtheorem{corollary}[theorem]{Corollary}

\theoremstyle{definition}
\newtheorem{definition}[theorem]{Definition}
\newtheorem{example}[theorem]{Example}
\newtheorem{construction}[theorem]{Construction}

\theoremstyle{remark}
\newtheorem{remark}[theorem]{Remark}
\newtheorem{notation}[theorem]{Notation}

% ── Python listings style ──────────────────────────────────────────────
\definecolor{jugeoBlue}{HTML}{2563EB}
\definecolor{jugeoGreen}{HTML}{059669}
\definecolor{jugeoRed}{HTML}{DC2626}
\definecolor{jugeoPurple}{HTML}{7C3AED}
\definecolor{jugeoGray}{HTML}{6B7280}
\definecolor{jugeoBg}{HTML}{F8FAFC}

\lstdefinestyle{jugeo-python}{
  language=Python,
  basicstyle=\ttfamily\small,
  keywordstyle=\color{jugeoBlue}\bfseries,
  stringstyle=\color{jugeoGreen},
  commentstyle=\color{jugeoGray}\itshape,
  numberstyle=\tiny\color{jugeoGray},
  backgroundcolor=\color{jugeoBg},
  numbers=left,
  numbersep=6pt,
  frame=single,
  framerule=0.4pt,
  rulecolor=\color{jugeoGray!40},
  breaklines=true,
  breakatwhitespace=true,
  showstringspaces=false,
  tabsize=4,
  xleftmargin=1.5em,
  framexleftmargin=1.5em,
  morekeywords={dataclass,frozen,field,Enum,IntEnum,auto,Any,
                Mapping,Sequence,Optional,match,case},
  literate={→}{{$\to$}}1 {←}{{$\leftarrow$}}1
           {≤}{{$\leq$}}1 {≥}{{$\geq$}}1 {≠}{{$\neq$}}1
           {∈}{{$\in$}}1 {∀}{{$\forall$}}1 {∃}{{$\exists$}}1
           {⊕}{{$\oplus$}}1 {⊖}{{$\ominus$}}1,
  escapeinside={(*@}{@*)},
}
\lstset{style=jugeo-python}

\lstdefinestyle{jugeo-output}{
  basicstyle=\ttfamily\small,
  backgroundcolor=\color{jugeoBg},
  frame=single,
  framerule=0.4pt,
  rulecolor=\color{jugeoGray!40},
  breaklines=true,
  showstringspaces=false,
  xleftmargin=1.5em,
  framexleftmargin=0.5em,
  numbers=none,
}

% ── JuGeo-specific macros ──────────────────────────────────────────────

% Judgment 8-tuple
\newcommand{\J}{\mathcal{J}}
\newcommand{\Jud}[8]{(#1,\,#2,\,#3,\,#4,\,#5,\,#6,\,#7,\,#8)}
\newcommand{\JudShort}[1]{\J_{#1}}

% Components
\newcommand{\coord}{\mathsf{c}}            % coordinate
\newcommand{\prop}{\varphi}                 % proposition
\newcommand{\carrier}{\mathcal{A}}          % carrier type
\newcommand{\evid}{\mathcal{E}}             % evidence bundle
\newcommand{\oblig}{\mathcal{O}}            % obligations
\newcommand{\obstr}{\mathcal{B}}            % obstructions
\newcommand{\trust}{\mathcal{T}}            % trust annotation
\newcommand{\prov}{\Pi}                     % provenance

% Trust algebra
\newcommand{\Talg}{\mathfrak{T}}            % trust ordered algebra
\newcommand{\Eadm}{\mathcal{E}_{\mathrm{adm}}}
\newcommand{\tleq}{\preceq}                 % trust ordering
\newcommand{\tjoin}{\oplus}                 % conservative join
\newcommand{\tatten}{\ominus}               % attenuation
\newcommand{\tpromote}{\uparrow_\pi}        % promotion
\newcommand{\tdemote}{\downarrow_\chi}       % demotion

% Trust tiers
\newcommand{\Tcontra}{\ensuremath{\mathsf{CONTRADICTED}}}
\newcommand{\Tunver}{\ensuremath{\mathsf{UNVERIFIED}}}
\newcommand{\Tcopilot}{\ensuremath{\mathsf{COPILOT}}}
\newcommand{\Toracle}{\ensuremath{\mathsf{ORACLE}}}
\newcommand{\Truntime}{\ensuremath{\mathsf{RUNTIME}}}
\newcommand{\Tsolver}{\ensuremath{\mathsf{SOLVER}}}
\newcommand{\Tproof}{\ensuremath{\mathsf{PROOF}}}

% Site and geometry
\newcommand{\Site}{\mathcal{S}}             % semantic site
\newcommand{\Cov}{\mathrm{Cov}}            % covering family
\newcommand{\Ob}{\mathrm{Ob}}              % objects
\newcommand{\Mor}{\mathrm{Mor}}            % morphisms
\newcommand{\res}{\mathrm{res}}            % restriction
\newcommand{\Cech}{\check{C}}              % Čech complex
\newcommand{\Hcoh}[1]{H^{#1}}             % cohomology group

% Descent
\newcommand{\descent}{\mathrm{desc}}
\newcommand{\glue}{\mathrm{glue}}
\newcommand{\overlap}[2]{#1 \cap #2}

% Semantic moves
\newcommand{\VERIFY}{\textsc{Verify}}
\newcommand{\CONSTRUCT}{\textsc{Construct}}
\newcommand{\NORMALIZE}{\textsc{Normalize}}
\newcommand{\GLUE}{\textsc{Glue}}
\newcommand{\RESTRICT}{\textsc{Restrict}}
\newcommand{\TRANSPORT}{\textsc{Transport}}
\newcommand{\REPAIR}{\textsc{Repair}}
\newcommand{\PROMOTE}{\textsc{Promote}}
\newcommand{\CHALLENGE}{\textsc{Challenge}}
\newcommand{\ESCALATE}{\textsc{Escalate}}

% Fragments
\newcommand{\QFLIA}{\texttt{QF\_LIA}}
\newcommand{\QFLRA}{\texttt{QF\_LRA}}
\newcommand{\QFBV}{\texttt{QF\_BV}}
\newcommand{\QFUF}{\texttt{QF\_UF}}

% Misc
\newcommand{\Python}{\textsc{Python}}
\newcommand{\JuGeo}{\textsc{JuGeo}}
\newcommand{\LEAN}{\textsc{Lean}\,4}
\newcommand{\Fstar}{F$^{\star}$}
\newcommand{\Coq}{\textsc{Coq}}
\newcommand{\Dafny}{\textsc{Dafny}}

% Common shortcuts
\newcommand{\ie}{i.e.\@\xspace}
\newcommand{\eg}{e.g.\@\xspace}
\newcommand{\cf}{cf.\@\xspace}
\newcommand{\wrt}{w.r.t.\@\xspace}

% Evaluation shortcuts
\newcommand{\numcases}{300}
\newcommand{\numspec}{100}
\newcommand{\numeq}{100}
\newcommand{\numbug}{100}
\newcommand{\medtime}{6.5\,ms}
\newcommand{\totaltime}{1.949\,s}


% ── Packages ────────────────────────────────────────────────────────────
\usepackage{amsmath,amssymb,amsthm,mathtools}
\usepackage{tikz-cd}
\usepackage{listings}
\usepackage{xcolor}
\usepackage{xspace}
\usepackage{booktabs}
\usepackage{multirow}
\usepackage{enumitem}
\usepackage[expansion=false]{microtype}
\usepackage{hyperref}
\usepackage{cleveref}
% \usepackage{thmtools}  % disabled: not available in this TeX installation
\usepackage{float}
\usepackage{caption}
\usepackage{subcaption}
\usepackage{tabularx}
\usepackage{stmaryrd}       % \llbracket, \rrbracket
\usepackage{bbm}            % \mathbbm{1}

% ── Page geometry ───────────────────────────────────────────────────────
\usepackage[margin=1in]{geometry}

% ── Theorem environments ────────────────────────────────────────────────
\theoremstyle{plain}
\newtheorem{theorem}{Theorem}[section]
\newtheorem{lemma}[theorem]{Lemma}
\newtheorem{proposition}[theorem]{Proposition}
\newtheorem{corollary}[theorem]{Corollary}

\theoremstyle{definition}
\newtheorem{definition}[theorem]{Definition}
\newtheorem{example}[theorem]{Example}
\newtheorem{construction}[theorem]{Construction}

\theoremstyle{remark}
\newtheorem{remark}[theorem]{Remark}
\newtheorem{notation}[theorem]{Notation}

% ── Python listings style ──────────────────────────────────────────────
\definecolor{jugeoBlue}{HTML}{2563EB}
\definecolor{jugeoGreen}{HTML}{059669}
\definecolor{jugeoRed}{HTML}{DC2626}
\definecolor{jugeoPurple}{HTML}{7C3AED}
\definecolor{jugeoGray}{HTML}{6B7280}
\definecolor{jugeoBg}{HTML}{F8FAFC}

\lstdefinestyle{jugeo-python}{
  language=Python,
  basicstyle=\ttfamily\small,
  keywordstyle=\color{jugeoBlue}\bfseries,
  stringstyle=\color{jugeoGreen},
  commentstyle=\color{jugeoGray}\itshape,
  numberstyle=\tiny\color{jugeoGray},
  backgroundcolor=\color{jugeoBg},
  numbers=left,
  numbersep=6pt,
  frame=single,
  framerule=0.4pt,
  rulecolor=\color{jugeoGray!40},
  breaklines=true,
  breakatwhitespace=true,
  showstringspaces=false,
  tabsize=4,
  xleftmargin=1.5em,
  framexleftmargin=1.5em,
  morekeywords={dataclass,frozen,field,Enum,IntEnum,auto,Any,
                Mapping,Sequence,Optional,match,case},
  literate={→}{{$\to$}}1 {←}{{$\leftarrow$}}1
           {≤}{{$\leq$}}1 {≥}{{$\geq$}}1 {≠}{{$\neq$}}1
           {∈}{{$\in$}}1 {∀}{{$\forall$}}1 {∃}{{$\exists$}}1
           {⊕}{{$\oplus$}}1 {⊖}{{$\ominus$}}1,
  escapeinside={(*@}{@*)},
}
\lstset{style=jugeo-python}

\lstdefinestyle{jugeo-output}{
  basicstyle=\ttfamily\small,
  backgroundcolor=\color{jugeoBg},
  frame=single,
  framerule=0.4pt,
  rulecolor=\color{jugeoGray!40},
  breaklines=true,
  showstringspaces=false,
  xleftmargin=1.5em,
  framexleftmargin=0.5em,
  numbers=none,
}

% ── JuGeo-specific macros ──────────────────────────────────────────────

% Judgment 8-tuple
\newcommand{\J}{\mathcal{J}}
\newcommand{\Jud}[8]{(#1,\,#2,\,#3,\,#4,\,#5,\,#6,\,#7,\,#8)}
\newcommand{\JudShort}[1]{\J_{#1}}

% Components
\newcommand{\coord}{\mathsf{c}}            % coordinate
\newcommand{\prop}{\varphi}                 % proposition
\newcommand{\carrier}{\mathcal{A}}          % carrier type
\newcommand{\evid}{\mathcal{E}}             % evidence bundle
\newcommand{\oblig}{\mathcal{O}}            % obligations
\newcommand{\obstr}{\mathcal{B}}            % obstructions
\newcommand{\trust}{\mathcal{T}}            % trust annotation
\newcommand{\prov}{\Pi}                     % provenance

% Trust algebra
\newcommand{\Talg}{\mathfrak{T}}            % trust ordered algebra
\newcommand{\Eadm}{\mathcal{E}_{\mathrm{adm}}}
\newcommand{\tleq}{\preceq}                 % trust ordering
\newcommand{\tjoin}{\oplus}                 % conservative join
\newcommand{\tatten}{\ominus}               % attenuation
\newcommand{\tpromote}{\uparrow_\pi}        % promotion
\newcommand{\tdemote}{\downarrow_\chi}       % demotion

% Trust tiers
\newcommand{\Tcontra}{\ensuremath{\mathsf{CONTRADICTED}}}
\newcommand{\Tunver}{\ensuremath{\mathsf{UNVERIFIED}}}
\newcommand{\Tcopilot}{\ensuremath{\mathsf{COPILOT}}}
\newcommand{\Toracle}{\ensuremath{\mathsf{ORACLE}}}
\newcommand{\Truntime}{\ensuremath{\mathsf{RUNTIME}}}
\newcommand{\Tsolver}{\ensuremath{\mathsf{SOLVER}}}
\newcommand{\Tproof}{\ensuremath{\mathsf{PROOF}}}

% Site and geometry
\newcommand{\Site}{\mathcal{S}}             % semantic site
\newcommand{\Cov}{\mathrm{Cov}}            % covering family
\newcommand{\Ob}{\mathrm{Ob}}              % objects
\newcommand{\Mor}{\mathrm{Mor}}            % morphisms
\newcommand{\res}{\mathrm{res}}            % restriction
\newcommand{\Cech}{\check{C}}              % Čech complex
\newcommand{\Hcoh}[1]{H^{#1}}             % cohomology group

% Descent
\newcommand{\descent}{\mathrm{desc}}
\newcommand{\glue}{\mathrm{glue}}
\newcommand{\overlap}[2]{#1 \cap #2}

% Semantic moves
\newcommand{\VERIFY}{\textsc{Verify}}
\newcommand{\CONSTRUCT}{\textsc{Construct}}
\newcommand{\NORMALIZE}{\textsc{Normalize}}
\newcommand{\GLUE}{\textsc{Glue}}
\newcommand{\RESTRICT}{\textsc{Restrict}}
\newcommand{\TRANSPORT}{\textsc{Transport}}
\newcommand{\REPAIR}{\textsc{Repair}}
\newcommand{\PROMOTE}{\textsc{Promote}}
\newcommand{\CHALLENGE}{\textsc{Challenge}}
\newcommand{\ESCALATE}{\textsc{Escalate}}

% Fragments
\newcommand{\QFLIA}{\texttt{QF\_LIA}}
\newcommand{\QFLRA}{\texttt{QF\_LRA}}
\newcommand{\QFBV}{\texttt{QF\_BV}}
\newcommand{\QFUF}{\texttt{QF\_UF}}

% Misc
\newcommand{\Python}{\textsc{Python}}
\newcommand{\JuGeo}{\textsc{JuGeo}}
\newcommand{\LEAN}{\textsc{Lean}\,4}
\newcommand{\Fstar}{F$^{\star}$}
\newcommand{\Coq}{\textsc{Coq}}
\newcommand{\Dafny}{\textsc{Dafny}}

% Common shortcuts
\newcommand{\ie}{i.e.\@\xspace}
\newcommand{\eg}{e.g.\@\xspace}
\newcommand{\cf}{cf.\@\xspace}
\newcommand{\wrt}{w.r.t.\@\xspace}

% Evaluation shortcuts
\newcommand{\numcases}{300}
\newcommand{\numspec}{100}
\newcommand{\numeq}{100}
\newcommand{\numbug}{100}
\newcommand{\medtime}{6.5\,ms}
\newcommand{\totaltime}{1.949\,s}

% experiment-data.tex — AUTO-GENERATED from real jugeo CLI runs
% DO NOT EDIT — regenerate with: python3 experiments/run_universal_metrics.py
% Generated from 30 programs across 6 domains

\newcommand{\expTotalPrograms}{30}
\newcommand{\expVerified}{30}
\newcommand{\expAccuracy}{100.0\%}
\newcommand{\expCoordsMin}{2}
\newcommand{\expCoordsMax}{10}
\newcommand{\expCoordsMean}{5.2}
\newcommand{\expCoordsMedian}{5.5}
\newcommand{\expPropsMin}{4}
\newcommand{\expPropsMax}{20}
\newcommand{\expPropsMean}{10.4}
\newcommand{\expPropsSum}{311}
\newcommand{\expPropsOkSum}{310}
\newcommand{\expTimeMin}{3.506\,s}
\newcommand{\expTimeMax}{6.714\,s}
\newcommand{\expTimeMean}{4.64\,s}
\newcommand{\expTimeMedian}{4.508\,s}
\newcommand{\expTimeTotal}{139.204\,s}
\newcommand{\expObstructionTotal}{0}
\newcommand{\expHOneZero}{30}
\newcommand{\expDomainCount}{6}
\newcommand{\expTrustCOPILOTSUGGESTED}{30}
\newcommand{\expDomainDsCount}{5}
\newcommand{\expDomainDsVerified}{5}
\newcommand{\expDomainDsAvgCoords}{7.6}
\newcommand{\expDomainDsAvgProps}{15.2}
\newcommand{\expDomainMathCount}{5}
\newcommand{\expDomainMathVerified}{5}
\newcommand{\expDomainMathAvgCoords}{4.8}
\newcommand{\expDomainMathAvgProps}{9.4}
\newcommand{\expDomainSmCount}{5}
\newcommand{\expDomainSmVerified}{5}
\newcommand{\expDomainSmAvgCoords}{7.6}
\newcommand{\expDomainSmAvgProps}{15.2}
\newcommand{\expDomainSortCount}{5}
\newcommand{\expDomainSortVerified}{5}
\newcommand{\expDomainSortAvgCoords}{2.4}
\newcommand{\expDomainSortAvgProps}{4.8}
\newcommand{\expDomainStrCount}{5}
\newcommand{\expDomainStrVerified}{5}
\newcommand{\expDomainStrAvgCoords}{3.2}
\newcommand{\expDomainStrAvgProps}{6.4}
\newcommand{\expDomainWebCount}{5}
\newcommand{\expDomainWebVerified}{5}
\newcommand{\expDomainWebAvgCoords}{5.6}
\newcommand{\expDomainWebAvgProps}{11.2}

% subsystem-data.tex — AUTO-GENERATED from real jugeo subsystem experiments
% DO NOT EDIT — regenerate with: python3 experiments/run_subsystem_experiments.py

\newcommand{\subCliTotal}{10}
\newcommand{\subCliVerified}{9}
\newcommand{\subCliAccuracy}{90.0\%}
\newcommand{\subCliMeanTime}{4.132\,s}
\newcommand{\subCliMinTime}{3.472\,s}
\newcommand{\subCliMaxTime}{5.372\,s}
\newcommand{\subCliTotalCoords}{54}
\newcommand{\subCliTotalProps}{107}
\newcommand{\subCliTotalPropsOk}{106}
\newcommand{\subSiteCalls}{90}
\newcommand{\subSiteSuccessful}{69}
\newcommand{\subSiteSuccessRate}{76.7\%}
\newcommand{\subSiteMeanTime}{0.0001\,s}
\newcommand{\subSiteTotalTime}{0.005\,s}
\newcommand{\subSpecificationsatisfactionSmallTime}{0.0003\,s}
\newcommand{\subSpecificationsatisfactionMediumTime}{0.0001\,s}
\newcommand{\subSpecificationsatisfactionLargeTime}{0.0001\,s}
\newcommand{\subInhabitantfleetSmallTime}{0.0\,s}
\newcommand{\subInhabitantfleetMediumTime}{0.0\,s}
\newcommand{\subInhabitantfleetLargeTime}{0.0\,s}
\newcommand{\subTheoremecologySmallTime}{0.0\,s}
\newcommand{\subTheoremecologyMediumTime}{0.0\,s}
\newcommand{\subTheoremecologyLargeTime}{0.0\,s}
\newcommand{\subStatespaceexplorationSmallTime}{0.0\,s}
\newcommand{\subStatespaceexplorationMediumTime}{0.0\,s}
\newcommand{\subStatespaceexplorationLargeTime}{0.0\,s}
\newcommand{\subRepairsemanticsSmallTime}{0.0\,s}
\newcommand{\subRepairsemanticsMediumTime}{0.0\,s}
\newcommand{\subRepairsemanticsLargeTime}{0.0\,s}
\newcommand{\subReplaygluingSmallTime}{0.0\,s}
\newcommand{\subReplaygluingMediumTime}{0.0\,s}
\newcommand{\subReplaygluingLargeTime}{0.0\,s}
\newcommand{\subSemanticfuturesSmallTime}{0.0\,s}
\newcommand{\subSemanticfuturesMediumTime}{0.0\,s}
\newcommand{\subSemanticfuturesLargeTime}{0.0\,s}
\newcommand{\subHypercovertreatySmallTime}{0.0\,s}
\newcommand{\subHypercovertreatyMediumTime}{0.0\,s}
\newcommand{\subHypercovertreatyLargeTime}{0.0\,s}
\newcommand{\subEncodeforsolverSmallTime}{0.0026\,s}
\newcommand{\subEncodeforsolverMediumTime}{0.0002\,s}
\newcommand{\subEncodeforsolverLargeTime}{0.0001\,s}
\newcommand{\subInterfaceroutingSmallTime}{0.0\,s}
\newcommand{\subInterfaceroutingMediumTime}{0.0\,s}
\newcommand{\subInterfaceroutingLargeTime}{0.0\,s}
\newcommand{\subOrchestrateverificationSmallTime}{0.0002\,s}
\newcommand{\subOrchestrateverificationMediumTime}{0.0\,s}
\newcommand{\subOrchestrateverificationLargeTime}{0.0\,s}
\newcommand{\subRunfulldescentSmallTime}{0.0\,s}
\newcommand{\subRunfulldescentMediumTime}{0.0\,s}
\newcommand{\subRunfulldescentLargeTime}{0.0\,s}
\newcommand{\subKernellifecycleSmallTime}{0.0\,s}
\newcommand{\subKernellifecycleMediumTime}{0.0\,s}
\newcommand{\subKernellifecycleLargeTime}{0.0\,s}
\newcommand{\subDiscoverypipelineSmallTime}{0.0\,s}
\newcommand{\subDiscoverypipelineMediumTime}{0.0\,s}
\newcommand{\subDiscoverypipelineLargeTime}{0.0\,s}
\newcommand{\subAnalogytransportSmallTime}{0.0\,s}
\newcommand{\subAnalogytransportMediumTime}{0.0\,s}
\newcommand{\subAnalogytransportLargeTime}{0.0\,s}
\newcommand{\subFormalcoresiteSmallTime}{0.0001\,s}
\newcommand{\subFormalcoresiteMediumTime}{0.0\,s}
\newcommand{\subFormalcoresiteLargeTime}{0.0\,s}
\newcommand{\subProblematlasSmallTime}{0.0\,s}
\newcommand{\subProblematlasMediumTime}{0.0\,s}
\newcommand{\subProblematlasLargeTime}{0.0\,s}
\newcommand{\subPublicalignmentSmallTime}{0.0\,s}
\newcommand{\subPublicalignmentMediumTime}{0.0\,s}
\newcommand{\subPublicalignmentLargeTime}{0.0\,s}
\newcommand{\subRegimebootstrappingSmallTime}{0.0\,s}
\newcommand{\subRegimebootstrappingMediumTime}{0.0\,s}
\newcommand{\subRegimebootstrappingLargeTime}{0.0\,s}
\newcommand{\subRelationalrefinementSmallTime}{0.0\,s}
\newcommand{\subRelationalrefinementMediumTime}{0.0\,s}
\newcommand{\subRelationalrefinementLargeTime}{0.0\,s}
\newcommand{\subSemanticclosureSmallTime}{0.0\,s}
\newcommand{\subSemanticclosureMediumTime}{0.0\,s}
\newcommand{\subSemanticclosureLargeTime}{0.0\,s}
\newcommand{\subTheoremeconomicsSmallTime}{0.0001\,s}
\newcommand{\subTheoremeconomicsMediumTime}{0.0\,s}
\newcommand{\subTheoremeconomicsLargeTime}{0.0\,s}
\newcommand{\subBenchmarksuiteSmallTime}{0.0\,s}
\newcommand{\subBenchmarksuiteMediumTime}{0.0\,s}
\newcommand{\subBenchmarksuiteLargeTime}{0.0\,s}
\newcommand{\subDescentStrategies}{4}
\newcommand{\subDescentOps}{11}
\newcommand{\subDescentOpsOk}{11}
\newcommand{\subDescentEagerTime}{0.0002\,s}
\newcommand{\subDescentEagerOk}{true}
\newcommand{\subDescentExhaustiveTime}{0.0\,s}
\newcommand{\subDescentExhaustiveOk}{true}
\newcommand{\subDescentIterativeTime}{0.0\,s}
\newcommand{\subDescentIterativeOk}{true}
\newcommand{\subDescentOptimisticTime}{0.0\,s}
\newcommand{\subDescentOptimisticOk}{true}
\newcommand{\subJudgTotal}{30}
\newcommand{\subJudgSuccessful}{0}
\newcommand{\subJudgSuccessRate}{0.0\%}
\newcommand{\subJudgMeanTimeUs}{7.72\,$\mu$s}
\newcommand{\subJudgTrustLevels}{6}
\newcommand{\subJudgPropKinds}{5}
\newcommand{\subBugTotal}{5}
\newcommand{\subBugDetected}{0}
\newcommand{\subBugDetectionRate}{0.0\%}
\newcommand{\subBugFalsePositiveRate}{0.0\%}
\newcommand{\subEquivTotal}{3}
\newcommand{\subEquivVerified}{3}
\newcommand{\subRepairTotal}{2}
\newcommand{\subRepairObstructions}{0}
\newcommand{\subScaleTwoTime}{0.0001\,s}
\newcommand{\subScaleFourTime}{0.0\,s}
\newcommand{\subScaleEightTime}{0.0\,s}
\newcommand{\subScaleTwelveTime}{0.0\,s}
\newcommand{\subScaleSixteenTime}{0.0\,s}
\newcommand{\subScaleTwentyTime}{0.0\,s}
\newcommand{\subTotalExperimentTime}{95.6\,s}

% comprehensive-data.tex — AUTO-GENERATED from real jugeo experiments
% DO NOT EDIT — regenerate with: python3 experiments/run_comprehensive_experiments.py
% Generated: 2026-03-23 09:19:06

% --- Size scaling metrics ---
\newcommand{\compTotalPrograms}{18}
\newcommand{\compTotalVerified}{18}
\newcommand{\compOverallAccuracy}{100.0\%}
\newcommand{\compTinyCount}{5}
\newcommand{\compTinyVerified}{5}
\newcommand{\compTinyMeanCoords}{3}
\newcommand{\compTinyMeanProps}{6}
\newcommand{\compTinyMeanTime}{4.95\,s}
\newcommand{\compTinyMeanLines}{5}
\newcommand{\compTinyMinTime}{4.45\,s}
\newcommand{\compTinyMaxTime}{5.51\,s}
\newcommand{\compSmallCount}{5}
\newcommand{\compSmallVerified}{5}
\newcommand{\compSmallMeanCoords}{3.6}
\newcommand{\compSmallMeanProps}{7.2}
\newcommand{\compSmallMeanTime}{4.85\,s}
\newcommand{\compSmallMeanLines}{16.0}
\newcommand{\compSmallMinTime}{4.30\,s}
\newcommand{\compSmallMaxTime}{5.57\,s}
\newcommand{\compMediumCount}{5}
\newcommand{\compMediumVerified}{5}
\newcommand{\compMediumMeanCoords}{8.6}
\newcommand{\compMediumMeanProps}{17.2}
\newcommand{\compMediumMeanTime}{4.47\,s}
\newcommand{\compMediumMeanLines}{45.0}
\newcommand{\compMediumMinTime}{4.26\,s}
\newcommand{\compMediumMaxTime}{4.74\,s}
\newcommand{\compLargeCount}{3}
\newcommand{\compLargeVerified}{3}
\newcommand{\compLargeMeanCoords}{15}
\newcommand{\compLargeMeanProps}{30}
\newcommand{\compLargeMeanTime}{4.31\,s}
\newcommand{\compLargeMeanLines}{79.0}
\newcommand{\compLargeMinTime}{4.27\,s}
\newcommand{\compLargeMaxTime}{4.38\,s}
% --- Aggregate timing ---
\newcommand{\compTimeMean}{4.68\,s}
\newcommand{\compTimeMedian}{4.50\,s}
\newcommand{\compTimeMin}{4.265\,s}
\newcommand{\compTimeMax}{5.571\,s}
\newcommand{\compTimeTotal}{84.3\,s}
\newcommand{\compTimeStdev}{0.45\,s}
% --- Aggregate structure ---
\newcommand{\compCoordsMean}{6.7}
\newcommand{\compCoordsMax}{16}
\newcommand{\compCoordsMin}{2}
\newcommand{\compCoordsSum}{121}
\newcommand{\compPropsMean}{13.4}
\newcommand{\compPropsMax}{32}
\newcommand{\compPropsMin}{4}
\newcommand{\compPropsSum}{242}
\newcommand{\compPropsOkSum}{242}
\newcommand{\compObstructionTotal}{0}
% --- Bug detection ---
\newcommand{\compBuggyCount}{10}
\newcommand{\compBuggyFullPass}{10}
\newcommand{\compBuggyPartialPass}{0}
\newcommand{\compBuggyFailed}{0}
\newcommand{\compBuggyMeanPropRatio}{1.00}
\newcommand{\compCorrectMeanPropRatio}{1.00}
\newcommand{\compBuggyMeanTime}{5.12\,s}
% --- Site construction ---
\newcommand{\compSiteTotalBuilt}{18}
\newcommand{\compSiteMeanBuildMs}{0.05}
\newcommand{\compSiteMaxBuildMs}{0.14}
\newcommand{\compSiteMeanCoords}{6.5}
\newcommand{\compSiteMeanMorphisms}{14.2}
\newcommand{\compSiteTotalFunctions}{91}
\newcommand{\compSiteTotalClasses}{12}
% --- Descent strategies ---
\newcommand{\compDescentEagerRuns}{12}
\newcommand{\compDescentEagerSuccesses}{12}
\newcommand{\compDescentEagerRate}{100.0\%}
\newcommand{\compDescentEagerMeanMs}{0.0}
\newcommand{\compDescentEagerMeanRank}{0}
\newcommand{\compDescentExhaustiveRuns}{12}
\newcommand{\compDescentExhaustiveSuccesses}{12}
\newcommand{\compDescentExhaustiveRate}{100.0\%}
\newcommand{\compDescentExhaustiveMeanMs}{0.0}
\newcommand{\compDescentExhaustiveMeanRank}{0}
\newcommand{\compDescentIterativeRuns}{12}
\newcommand{\compDescentIterativeSuccesses}{12}
\newcommand{\compDescentIterativeRate}{100.0\%}
\newcommand{\compDescentIterativeMeanMs}{0.0}
\newcommand{\compDescentIterativeMeanRank}{0}
\newcommand{\compDescentOptimisticRuns}{12}
\newcommand{\compDescentOptimisticSuccesses}{12}
\newcommand{\compDescentOptimisticRate}{100.0\%}
\newcommand{\compDescentOptimisticMeanMs}{0.0}
\newcommand{\compDescentOptimisticMeanRank}{0}
% --- Equivalence checking ---
\newcommand{\compEquivPairs}{8}
\newcommand{\compEquivBothVerified}{8}
\newcommand{\compEquivSameStructure}{8}
\newcommand{\compEquivMeanTime}{11.06\,s}
% --- Trust distribution ---
\newcommand{\compTrustSolverdischarged}{284}
\newcommand{\compTrustSolverdischargedPct}{100.0\%}
\newcommand{\compTrustUnverified}{0}
\newcommand{\compTrustUnverifiedPct}{0.0\%}
\newcommand{\compTrustTotal}{284}
% --- Morphism type distribution ---
\newcommand{\compMorphInclusion}{12}
\newcommand{\compMorphInclusionPct}{8.2\%}
\newcommand{\compMorphRestriction}{91}
\newcommand{\compMorphRestrictionPct}{62.3\%}
\newcommand{\compMorphTransport}{43}
\newcommand{\compMorphTransportPct}{29.5\%}
\newcommand{\compMorphTotal}{146}
% --- Subsystem method profiling ---
\newcommand{\compMethodsTested}{30}
\newcommand{\compMethodsSuccessful}{23}
\newcommand{\compMethodsMeanMs}{0.02}
\newcommand{\compProfAnalogyTransportMeanMs}{0.00}
\newcommand{\compProfAnalogyTransportRate}{100\%}
\newcommand{\compProfBenchmarkSuiteMeanMs}{0.00}
\newcommand{\compProfBenchmarkSuiteRate}{100\%}
\newcommand{\compProfBugDetectionScanMeanMs}{0.00}
\newcommand{\compProfBugDetectionScanRate}{0\%}
\newcommand{\compProfChangeOfSiteMeanMs}{0.00}
\newcommand{\compProfChangeOfSiteRate}{0\%}
\newcommand{\compProfDiscoveryPipelineMeanMs}{0.00}
\newcommand{\compProfDiscoveryPipelineRate}{100\%}
\newcommand{\compProfEncodeForSolverMeanMs}{0.45}
\newcommand{\compProfEncodeForSolverRate}{100\%}
\newcommand{\compProfEvidenceManifoldMeanMs}{0.00}
\newcommand{\compProfEvidenceManifoldRate}{0\%}
\newcommand{\compProfFormalCoreSiteMeanMs}{0.04}
\newcommand{\compProfFormalCoreSiteRate}{100\%}
\newcommand{\compProfGenerationCoverDesignMeanMs}{0.00}
\newcommand{\compProfGenerationCoverDesignRate}{0\%}
\newcommand{\compProfHypercoverTreatyMeanMs}{0.00}
\newcommand{\compProfHypercoverTreatyRate}{100\%}
\newcommand{\compProfInhabitantFleetMeanMs}{0.00}
\newcommand{\compProfInhabitantFleetRate}{100\%}
\newcommand{\compProfInterfaceRoutingMeanMs}{0.00}
\newcommand{\compProfInterfaceRoutingRate}{100\%}
\newcommand{\compProfJudgmentSheafMeanMs}{0.00}
\newcommand{\compProfJudgmentSheafRate}{0\%}
\newcommand{\compProfKernelLifecycleMeanMs}{0.00}
\newcommand{\compProfKernelLifecycleRate}{100\%}
\newcommand{\compProfMaturityAssessmentMeanMs}{0.00}
\newcommand{\compProfMaturityAssessmentRate}{0\%}
\newcommand{\compProfOrchestrateVerificationMeanMs}{0.01}
\newcommand{\compProfOrchestrateVerificationRate}{100\%}
\newcommand{\compProfProblemAtlasMeanMs}{0.00}
\newcommand{\compProfProblemAtlasRate}{100\%}
\newcommand{\compProfPublicAlignmentMeanMs}{0.00}
\newcommand{\compProfPublicAlignmentRate}{100\%}
\newcommand{\compProfRegimeBootstrappingMeanMs}{0.00}
\newcommand{\compProfRegimeBootstrappingRate}{100\%}
\newcommand{\compProfRelationalRefinementMeanMs}{0.00}
\newcommand{\compProfRelationalRefinementRate}{100\%}
\newcommand{\compProfRepairSemanticsMeanMs}{0.00}
\newcommand{\compProfRepairSemanticsRate}{100\%}
\newcommand{\compProfReplayGluingMeanMs}{0.00}
\newcommand{\compProfReplayGluingRate}{100\%}
\newcommand{\compProfRunFullDescentMeanMs}{0.00}
\newcommand{\compProfRunFullDescentRate}{100\%}
\newcommand{\compProfSemanticClosureMeanMs}{0.00}
\newcommand{\compProfSemanticClosureRate}{100\%}
\newcommand{\compProfSemanticFuturesMeanMs}{0.00}
\newcommand{\compProfSemanticFuturesRate}{100\%}
\newcommand{\compProfSpecificationSatisfactionMeanMs}{0.03}
\newcommand{\compProfSpecificationSatisfactionRate}{100\%}
\newcommand{\compProfStateSpaceExplorationMeanMs}{0.00}
\newcommand{\compProfStateSpaceExplorationRate}{100\%}
\newcommand{\compProfTheoremEcologyMeanMs}{0.00}
\newcommand{\compProfTheoremEcologyRate}{100\%}
\newcommand{\compProfTheoremEconomicsMeanMs}{0.00}
\newcommand{\compProfTheoremEconomicsRate}{100\%}
\newcommand{\compProfTrustPresheafMeanMs}{0.00}
\newcommand{\compProfTrustPresheafRate}{0\%}

% data-paper68.tex — AUTO-GENERATED by exp68_technical_debt.py
% DO NOT EDIT — regenerate with: python3 experiments/exp68_technical_debt.py
% Generated from 10 programs (clean + debt-laden)

\newcommand{\ppLXVIIIprogramCount}{10}
\newcommand{\ppLXVIIIprogramsOk}{10}
\newcommand{\ppLXVIIIverified}{10}
\newcommand{\ppLXVIIIverifiedPct}{100.0\%}
\newcommand{\ppLXVIIIcleanCount}{5}
\newcommand{\ppLXVIIIdebtCount}{5}

\newcommand{\ppLXVIIIcoordsMean}{5.5}
\newcommand{\ppLXVIIIcoordsMax}{9}
\newcommand{\ppLXVIIIcoordsMin}{3}
\newcommand{\ppLXVIIIpropsMean}{11}
\newcommand{\ppLXVIIIpropsMax}{18}
\newcommand{\ppLXVIIIpropsMin}{6}
\newcommand{\ppLXVIIIpropsTotal}{110}
\newcommand{\ppLXVIIIbugsMean}{0}
\newcommand{\ppLXVIIIbugsTotal}{0}

\newcommand{\ppLXVIIIcleanCoordsMean}{5.6}
\newcommand{\ppLXVIIIcleanPropsMean}{11.2}
\newcommand{\ppLXVIIIcleanTimeMean}{19.85\,s}
\newcommand{\ppLXVIIIdebtCoordsMean}{5.4}
\newcommand{\ppLXVIIIdebtPropsMean}{10.8}
\newcommand{\ppLXVIIIdebtTimeMean}{18.35\,s}

\newcommand{\ppLXVIIIcycleMean}{50.0\,ms}
\newcommand{\ppLXVIIItimeMean}{19.10\,s}
\newcommand{\ppLXVIIItimeTotal}{190.96\,s}
\newcommand{\ppLXVIIItimeMax}{28.739\,s}
\newcommand{\ppLXVIIItimeMin}{14.669\,s}

% Per-program debt results
\newcommand{\ppLXVIIIdebtcleanstackScore}{0.5}
\newcommand{\ppLXVIIIdebtcleanstackTrust}{0.5}
\newcommand{\ppLXVIIIdebtcleanstackCoords}{8}
\newcommand{\ppLXVIIIdebtcleanstackProps}{16}
\newcommand{\ppLXVIIIdebtcleanstackBugs}{0}
\newcommand{\ppLXVIIIdebtcleanstackTime}{17.628\,s}
\newcommand{\ppLXVIIIdebtcleancalculatorScore}{0.5}
\newcommand{\ppLXVIIIdebtcleancalculatorTrust}{0.5}
\newcommand{\ppLXVIIIdebtcleancalculatorCoords}{7}
\newcommand{\ppLXVIIIdebtcleancalculatorProps}{14}
\newcommand{\ppLXVIIIdebtcleancalculatorBugs}{0}
\newcommand{\ppLXVIIIdebtcleancalculatorTime}{17.65\,s}
\newcommand{\ppLXVIIIdebtmoderatedebtparserScore}{0.5}
\newcommand{\ppLXVIIIdebtmoderatedebtparserTrust}{0.5}
\newcommand{\ppLXVIIIdebtmoderatedebtparserCoords}{4}
\newcommand{\ppLXVIIIdebtmoderatedebtparserProps}{8}
\newcommand{\ppLXVIIIdebtmoderatedebtparserBugs}{0}
\newcommand{\ppLXVIIIdebtmoderatedebtparserTime}{16.294\,s}
\newcommand{\ppLXVIIIdebthighdebtglobalsScore}{0.5}
\newcommand{\ppLXVIIIdebthighdebtglobalsTrust}{0.5}
\newcommand{\ppLXVIIIdebthighdebtglobalsCoords}{4}
\newcommand{\ppLXVIIIdebthighdebtglobalsProps}{8}
\newcommand{\ppLXVIIIdebthighdebtglobalsBugs}{0}
\newcommand{\ppLXVIIIdebthighdebtglobalsTime}{17.528\,s}
\newcommand{\ppLXVIIIdebtdebtmixedconcernsScore}{0.5}
\newcommand{\ppLXVIIIdebtdebtmixedconcernsTrust}{0.5}
\newcommand{\ppLXVIIIdebtdebtmixedconcernsCoords}{6}
\newcommand{\ppLXVIIIdebtdebtmixedconcernsProps}{12}
\newcommand{\ppLXVIIIdebtdebtmixedconcernsBugs}{0}
\newcommand{\ppLXVIIIdebtdebtmixedconcernsTime}{17.484\,s}
\newcommand{\ppLXVIIIdebtcleanvalidatorScore}{0.5}
\newcommand{\ppLXVIIIdebtcleanvalidatorTrust}{0.5}
\newcommand{\ppLXVIIIdebtcleanvalidatorCoords}{4}
\newcommand{\ppLXVIIIdebtcleanvalidatorProps}{8}
\newcommand{\ppLXVIIIdebtcleanvalidatorBugs}{0}
\newcommand{\ppLXVIIIdebtcleanvalidatorTime}{16.233\,s}
\newcommand{\ppLXVIIIdebtdebtnoerrorhandlingScore}{0.5}
\newcommand{\ppLXVIIIdebtdebtnoerrorhandlingTrust}{0.5}
\newcommand{\ppLXVIIIdebtdebtnoerrorhandlingCoords}{4}
\newcommand{\ppLXVIIIdebtdebtnoerrorhandlingProps}{8}
\newcommand{\ppLXVIIIdebtdebtnoerrorhandlingBugs}{0}
\newcommand{\ppLXVIIIdebtdebtnoerrorhandlingTime}{14.669\,s}
\newcommand{\ppLXVIIIdebtcleansortingScore}{0.5}
\newcommand{\ppLXVIIIdebtcleansortingTrust}{0.5}
\newcommand{\ppLXVIIIdebtcleansortingCoords}{3}
\newcommand{\ppLXVIIIdebtcleansortingProps}{6}
\newcommand{\ppLXVIIIdebtcleansortingBugs}{0}
\newcommand{\ppLXVIIIdebtcleansortingTime}{18.975\,s}
\newcommand{\ppLXVIIIdebtdebtgodclassScore}{0.5}
\newcommand{\ppLXVIIIdebtdebtgodclassTrust}{0.5}
\newcommand{\ppLXVIIIdebtdebtgodclassCoords}{9}
\newcommand{\ppLXVIIIdebtdebtgodclassProps}{18}
\newcommand{\ppLXVIIIdebtdebtgodclassBugs}{0}
\newcommand{\ppLXVIIIdebtdebtgodclassTime}{25.758\,s}
\newcommand{\ppLXVIIIdebtcleantreeScore}{0.5}
\newcommand{\ppLXVIIIdebtcleantreeTrust}{0.5}
\newcommand{\ppLXVIIIdebtcleantreeCoords}{6}
\newcommand{\ppLXVIIIdebtcleantreeProps}{12}
\newcommand{\ppLXVIIIdebtcleantreeBugs}{0}
\newcommand{\ppLXVIIIdebtcleantreeTime}{28.739\,s}


% ── Additional packages ────────────────────────────────────────────
\usepackage{array}
\usepackage{tikz}
\usetikzlibrary{arrows.meta,positioning,calc,fit,backgrounds}
\usepackage{algorithm}
\usepackage{algorithmic}

% ── Paper-specific macros ──────────────────────────────────────────
\newcommand{\TechDebt}{\textsc{TechDebt}}
\newcommand{\MatAssess}{\textsc{MaturityAssessor}}
\newcommand{\DebtScore}{\mathsf{debt\_score}}
\newcommand{\TrustDecay}{\mathsf{trust\_decay}}
\newcommand{\MaturityLevel}{\mathsf{maturity\_level}}
\newcommand{\MaturityAssessment}{\texttt{maturity\_assessment}}
\newcommand{\TrustPresheaf}{\texttt{trust\_presheaf}}
\newcommand{\OrchestrateVerif}{\texttt{orchestrate\_verification}}
\newcommand{\EncodeForSolver}{\texttt{encode\_for\_solver}}
\newcommand{\CopilotSugg}{\mathsf{COPILOT\_SUGGESTED}}
\newcommand{\SolverDisch}{\mathsf{SOLVER\_DISCHARGED}}
\newcommand{\TrustUpgrade}{\rightsquigarrow}
\newcommand{\DebtThreshold}{\delta_{\mathrm{thresh}}}

\title{\textbf{Measuring and Managing Technical Debt\\
via Trust Degradation and Maturity Assessment}}

\author{JuGeo Research Group}
\date{\today}

\begin{document}
\maketitle

% ─────────────────────────────────────────────────────────────────────
\begin{abstract}
Technical debt accumulates when software evolves without adequate
verification, causing trust in correctness guarantees to erode.
We present a sheaf-theoretic framework for measuring and managing
technical debt within \JuGeo\ by modeling trust degradation as a
presheaf morphism and maturity progression as a functor over the
semantic site. We introduce a debt score $\DebtScore(c)$ quantifying the gap
between current trust and target maturity, and we discuss how trust-presheaf
and maturity-style summaries could support repair planning. The Paper~68 data
are best understood as a prototype runtime characterization: across
\ppLXVIIIprogramCount{} curated programs (\ppLXVIIIcleanCount{} clean and
\ppLXVIIIdebtCount{} debt-laden), all instances verified, no bugs were
reported, and the generated per-program debt scores were uniformly 0.5.
Accordingly, this revision does not claim that the present implementation
successfully separates clean code from debt-laden code.
\end{abstract}

\newpage

% =====================================================================
\section{Introduction}\label{sec:intro}
% =====================================================================

Technical debt is pervasive in software engineering.  Code that was
once verified may lose correctness guarantees as dependencies change or
specifications evolve.  Traditional debt metrics---cyclomatic
complexity, coverage, lint warnings---correlate weakly with actual
correctness risk.

\paragraph{The trust perspective.}
\JuGeo\ offers a principled alternative: every coordinate $\coord$
carries a trust annotation $\trust \in \Talg$.  When code evolves,
trust may \emph{degrade}: a coordinate at $\Tsolver$ may fall to
$\Tcopilot$ or $\Tunver$ if assumptions are invalidated.  This trust
degradation is the formal manifestation of technical debt.

\paragraph{Contributions.}
\begin{itemize}[noitemsep]
  \item A formal model of technical debt as trust-maturity gap
        (§\ref{sec:debt-model}).
  \item The trust degradation presheaf (§\ref{sec:trust-degradation}).
  \item The \MaturityAssessment\ algorithm (§\ref{sec:maturity}).
  \item Debt-aware repair with convergence guarantees
        (§\ref{sec:repair}).
  \item Lean~4 proof sketches (§\ref{sec:lean-proof}).
  \item A prototype evaluation on \ppLXVIIIprogramCount{} curated programs
        showing runtime characteristics and current limitations
        (§\ref{sec:evaluation}).
\end{itemize}

% =====================================================================
\section{Background}\label{sec:background}
% =====================================================================

\JuGeo\ represents program analysis as sheaf theory over a semantic
site $\Site$.  Each coordinate $\coord$ carries a judgment
$\J = (\coord, \varphi, P, Q, I, \delta, \tau, \pi)$,
where $\tau \in \Talg$ orders as $\Tcontra \tleq \Tunver \tleq
\Tcopilot \tleq \Toracle \tleq \Truntime \tleq \Tsolver \tleq
\Tproof$.

Cunningham~\cite{Cunningham1992} introduced the debt metaphor.
Subsequent work~\cite{Kruchten2012,Ernst2015} classified debt
categories.  Quantitative approaches include SonarQube's SQALE
model~\cite{Letouzey2012}.  None ground debt in formal verification.

% =====================================================================
\section{Technical Debt Model}\label{sec:debt-model}
% =====================================================================

\begin{definition}[Target Maturity]\label{def:target}
A \emph{target maturity assignment} $\mu^* : \Ob(\Site) \to \{1,\ldots,5\}$
specifies the desired maturity for each coordinate.
\end{definition}

\begin{definition}[Debt Score]\label{def:debt-score}
For coordinate $\coord$ with current maturity $\mu(\coord)$ and target
$\mu^*(\coord)$, the debt score is
$\DebtScore(\coord) = \max(0, \mu^*(\coord) - \mu(\coord))$.
The aggregate debt is $D(\Site) = \sum_{\coord} w(\coord) \cdot
\DebtScore(\coord)$ where $w(\coord)$ is criticality weight.
\end{definition}

The five maturity levels map to minimum trust tiers:
\begin{center}
\begin{tabular}{lll}
\toprule
\textbf{Level} & \textbf{Min Trust} & \textbf{Description} \\
\midrule
\textsf{Initial}    & $\Tunver$   & No verification \\
\textsf{Managed}    & $\Tcopilot$ & LLM-suggested \\
\textsf{Defined}    & $\Toracle$  & Runtime-tested \\
\textsf{Measured}   & $\Tsolver$  & Solver-discharged \\
\textsf{Optimizing} & $\Tproof$   & Formally proved \\
\bottomrule
\end{tabular}
\end{center}

\begin{proposition}[Debt Monotonicity]\label{prop:debt-mono}
If $f : \coord_1 \to \coord_2$ is a morphism in $\Site$ and $\mu^*$
is monotone along $f$, then debt is compatible with the site structure:
restriction cannot artificially inflate debt in refinements.
\end{proposition}

% =====================================================================
\section{The Trust-Maturity Lattice}\label{sec:trust-maturity-lattice}
% =====================================================================

The five-stage maturity model introduced in \S\ref{sec:debt-model}
admits a rich algebraic structure.  In this section we develop the
\emph{trust-maturity lattice} as a bounded distributive lattice and
establish the embedding from $\Talg$ into it.

% ---------------------------------------------------------------------
\subsection{Lattice Axioms}\label{subsec:lattice-axioms}

\begin{definition}[Maturity Lattice]\label{def:maturity-lattice}
Let $\mathcal{M} = \{\textsf{Initial}, \textsf{Managed},
\textsf{Defined}, \textsf{Measured}, \textsf{Optimizing}\}$ be
equipped with the total order
\[
  \textsf{Initial} \sqsubseteq \textsf{Managed} \sqsubseteq
  \textsf{Defined} \sqsubseteq \textsf{Measured} \sqsubseteq
  \textsf{Optimizing}.
\]
Define binary operations $\sqcap$ (meet) and $\sqcup$ (join) as
\[
  a \sqcap b = \min(a,b), \qquad a \sqcup b = \max(a,b).
\]
Then $(\mathcal{M}, \sqsubseteq, \sqcap, \sqcup)$ is a lattice.
\end{definition}

\begin{theorem}[Bounded Distributive Lattice]\label{thm:bounded-dist}
$(\mathcal{M}, \sqsubseteq, \sqcap, \sqcup, \textsf{Initial},
\textsf{Optimizing})$ is a bounded distributive lattice.
\end{theorem}

\begin{proof}
We verify the lattice axioms:
\begin{enumerate}[noitemsep]
  \item \textbf{Idempotency:} $a \sqcap a = a$ and $a \sqcup a = a$
        follow from $\min(a,a)=a$ and $\max(a,a)=a$.
  \item \textbf{Commutativity:} $a \sqcap b = b \sqcap a$ and
        $a \sqcup b = b \sqcup a$ by commutativity of $\min$ and
        $\max$.
  \item \textbf{Associativity:} $a \sqcap (b \sqcap c) =
        (a \sqcap b) \sqcap c$ and similarly for $\sqcup$ by
        associativity of $\min$ and $\max$.
  \item \textbf{Absorption:} $a \sqcap (a \sqcup b) = a$ and
        $a \sqcup (a \sqcap b) = a$ since $\min(a,\max(a,b))=a$
        and $\max(a,\min(a,b))=a$.
  \item \textbf{Bounds:} $\textsf{Initial} \sqcap a = \textsf{Initial}$
        and $\textsf{Optimizing} \sqcup a = \textsf{Optimizing}$ for
        all $a$.
  \item \textbf{Distributivity:} For a totally ordered finite set,
        $a \sqcap (b \sqcup c) = (a \sqcap b) \sqcup (a \sqcap c)$
        holds by the identity
        $\min(a,\max(b,c)) = \max(\min(a,b),\min(a,c))$.
\end{enumerate}
\end{proof}

\begin{corollary}\label{cor:lattice-hom}
The debt score $\DebtScore : \mathcal{M} \times \mathcal{M} \to
\mathbb{N}$ is anti-monotone in the first argument and monotone
in the second: if $\mu_1 \sqsubseteq \mu_2$ then
$\DebtScore(\mu_2, \mu^*) \leq \DebtScore(\mu_1, \mu^*)$.
\end{corollary}

% ---------------------------------------------------------------------
\subsection{The Trust-to-Maturity Embedding}\label{subsec:trust-embed}

\begin{definition}[Trust-to-Maturity Map]\label{def:trust-to-mat}
Define $\iota : \Talg \to \mathcal{M}$ by:
\[
  \iota(\tau) = \begin{cases}
    \textsf{Initial}    & \text{if } \tau \tleq \Tunver, \\
    \textsf{Managed}    & \text{if } \tau = \Tcopilot, \\
    \textsf{Defined}    & \text{if } \tau \in \{\Toracle, \Truntime\}, \\
    \textsf{Measured}   & \text{if } \tau = \Tsolver, \\
    \textsf{Optimizing} & \text{if } \tau = \Tproof.
  \end{cases}
\]
\end{definition}

\begin{lemma}[Order-Preserving Embedding]\label{lem:order-pres}
The map $\iota$ is order-preserving: if $\tau_1 \tleq \tau_2$ then
$\iota(\tau_1) \sqsubseteq \iota(\tau_2)$.
\end{lemma}

\begin{proof}
By case analysis on the seven trust levels.  The trust ordering
$\Tcontra \tleq \Tunver \tleq \Tcopilot \tleq \Toracle \tleq
\Truntime \tleq \Tsolver \tleq \Tproof$ maps into the maturity
ordering via $\iota$, and one verifies that each consecutive pair
in $\Talg$ maps to a non-descending pair in $\mathcal{M}$.
\end{proof}

\begin{lemma}[Join-Compatibility]\label{lem:join-compat}
For all $\tau_1, \tau_2 \in \Talg$:
$\iota(\tau_1 \tjoin \tau_2) \sqsubseteq
\iota(\tau_1) \sqcup \iota(\tau_2)$.
\end{lemma}

\begin{proof}
Since $\tjoin$ is the join in $\Talg$ and $\iota$ is order-preserving
(Lemma~\ref{lem:order-pres}), we have
$\iota(\tau_1) \sqsubseteq \iota(\tau_1 \tjoin \tau_2)$ and
$\iota(\tau_2) \sqsubseteq \iota(\tau_1 \tjoin \tau_2)$, so
$\iota(\tau_1) \sqcup \iota(\tau_2) \sqsubseteq
\iota(\tau_1 \tjoin \tau_2)$.  Conversely, the coarsening from seven
to five levels can only collapse elements, giving the stated inequality.
\end{proof}

% ---------------------------------------------------------------------
\subsection{Lattice Diagram}\label{subsec:lattice-diagram}

\begin{figure}[H]
\centering
\begin{tikzpicture}[
  level/.style={draw, rounded corners=3pt, minimum width=2.2cm,
                minimum height=0.7cm, font=\small\sffamily},
  >=Stealth,
  node distance=1.0cm
]
  \node[level, fill=red!15]     (L1) {\textsf{Initial}};
  \node[level, fill=orange!15, above=of L1] (L2) {\textsf{Managed}};
  \node[level, fill=yellow!15, above=of L2] (L3) {\textsf{Defined}};
  \node[level, fill=green!15,  above=of L3] (L4) {\textsf{Measured}};
  \node[level, fill=blue!15,   above=of L4] (L5) {\textsf{Optimizing}};

  \draw[->, thick] (L1) -- (L2) node[midway, right, font=\footnotesize] {$\tpromote$};
  \draw[->, thick] (L2) -- (L3) node[midway, right, font=\footnotesize] {$\tpromote$};
  \draw[->, thick] (L3) -- (L4) node[midway, right, font=\footnotesize] {$\tpromote$};
  \draw[->, thick] (L4) -- (L5) node[midway, right, font=\footnotesize] {$\tpromote$};

  \draw[->, thick, dashed, jugeoRed] (L5) to[bend right=50]
    node[midway, left, font=\footnotesize, text=jugeoRed] {$\tdemote$} (L1);

  % Trust level annotations on the right
  \node[right=1.5cm of L1, font=\footnotesize, text=jugeoGray]
    {$\Tunver$, $\Tcontra$};
  \node[right=1.5cm of L2, font=\footnotesize, text=jugeoGray]
    {$\Tcopilot$};
  \node[right=1.5cm of L3, font=\footnotesize, text=jugeoGray]
    {$\Toracle$, $\Truntime$};
  \node[right=1.5cm of L4, font=\footnotesize, text=jugeoGray]
    {$\Tsolver$};
  \node[right=1.5cm of L5, font=\footnotesize, text=jugeoGray]
    {$\Tproof$};
\end{tikzpicture}
\caption{The five-stage maturity lattice $\mathcal{M}$ with
  trust-level annotations.  Solid arrows denote promotion
  ($\tpromote$); dashed arrow denotes degradation ($\tdemote$).}
\label{fig:maturity-lattice}
\end{figure}

\begin{definition}[Debt Gradient]\label{def:debt-gradient}
For a site $\Site$ with maturity assignment $\mu$ and target $\mu^*$,
the \emph{debt gradient} is the vector
\[
  \nabla D = \bigl(\DebtScore(\coord_1), \ldots,
  \DebtScore(\coord_n)\bigr) \in \mathbb{N}^n.
\]
The $\ell^1$-norm $\|\nabla D\|_1 = D(\Site)$ recovers the aggregate
debt.  The $\ell^\infty$-norm $\|\nabla D\|_\infty = \max_i
\DebtScore(\coord_i)$ identifies the worst-debt coordinate.
\end{definition}

\begin{lemma}[Repair Complexity Lower Bound]\label{lem:repair-lower}
Any repair strategy that reduces aggregate debt from $D$ to $0$ must
execute at least $\|\nabla D\|_1$ atomic promotion actions, where each
atomic action promotes one coordinate by exactly one maturity level.
\end{lemma}

\begin{proof}
Each atomic promotion reduces exactly one component of $\nabla D$ by
one, hence decreases $\|\nabla D\|_1$ by exactly one.  Since
$\|\nabla D\|_1$ must reach zero, the result follows.
\end{proof}



% =====================================================================
\section{Trust Degradation}\label{sec:trust-degradation}
% =====================================================================

Trust degrades via four triggers: dependency change, specification
drift, environment shift, and temporal decay.

\begin{definition}[Trust Degradation Presheaf]\label{def:degrade}
$\mathcal{D} : \Site^{\mathrm{op}} \to \mathbf{Set}$ assigns to each
$\coord$ the set of degradation events
$\{(t, \tau_{\mathrm{old}}, \tau_{\mathrm{new}}, \mathrm{reason})
\mid \tau_{\mathrm{new}} \tleq \tau_{\mathrm{old}}\}$.
Restriction maps propagate degradation along dependencies.
\end{definition}

\begin{theorem}[Degradation Propagation]\label{thm:degrade-prop}
Let $\coord_1 \to \coord_2$ be a dependency edge and
$d \in \mathcal{D}(\coord_2)$ a degradation event.  Then:
$\tau_{\mathrm{new}}(\coord_1) \tleq \tau_{\mathrm{old}}(\coord_1)
\tjoin \tau_{\mathrm{new}}(\coord_2)$.
\end{theorem}

\begin{proof}
By monotonicity of $\tjoin$ and the covering axiom of $\Site$.
\end{proof}

We model temporal decay as $\TrustDecay(\coord, t) = \tau_0(\coord)
\cdot e^{-\lambda_\coord t}$, where $\lambda_\coord$ is the decay
rate.  When $\TrustDecay$ drops below $\DebtThreshold$, the
coordinate is flagged for re-verification.

% =====================================================================
\section{Temporal Debt Evolution}\label{sec:temporal-evolution}
% =====================================================================

While \S\ref{sec:trust-degradation} introduced the exponential decay
model $\TrustDecay(\coord, t)$, we now develop a richer dynamical
systems perspective on how debt evolves over time.

\subsection{The Debt Velocity}\label{subsec:debt-velocity}

\begin{definition}[Debt Velocity]\label{def:debt-velocity}
Let $D(t)$ denote the aggregate debt at time $t$.  The \emph{debt
velocity} is
\[
  \dot{D}(t) = \frac{dD}{dt} = \sum_{\coord \in \Ob(\Site)}
  w(\coord) \cdot \frac{d}{dt}\DebtScore(\coord, t).
\]
Positive velocity indicates debt accumulation; negative velocity
indicates debt reduction.
\end{definition}

\begin{definition}[Decay Rate Matrix]\label{def:decay-matrix}
Let $\Lambda = \mathrm{diag}(\lambda_{\coord_1}, \ldots,
\lambda_{\coord_n})$ be the diagonal matrix of per-coordinate
decay rates.  The \emph{trust vector} $\mathbf{T}(t) =
(\tau(\coord_1, t), \ldots, \tau(\coord_n, t))^\top$ evolves as
\[
  \dot{\mathbf{T}}(t) = -\Lambda \, \mathbf{T}(t) + \mathbf{R}(t),
\]
where $\mathbf{R}(t)$ is the repair input vector encoding
verification actions taken at time $t$.
\end{definition}

\begin{theorem}[Super-Linear Accumulation]\label{thm:superlinear}
Without intervention ($\mathbf{R}(t) = \mathbf{0}$), aggregate debt
grows super-linearly.  Specifically, if all coordinates begin at
maturity $\mu_0$ with target $\mu^*$, then for $t$ large enough,
\[
  D(t) \geq \sum_{\coord} w(\coord) \cdot
  \bigl(\mu^*(\coord) - \mu_0 \cdot e^{-\lambda_\coord t}\bigr).
\]
As $t \to \infty$, $D(t) \to \sum_\coord w(\coord) \cdot \mu^*(\coord)$,
the maximum possible debt.
\end{theorem}

\begin{proof}
Without repair, $\tau(\coord, t) = \tau_0(\coord) \cdot
e^{-\lambda_\coord t}$, which is monotonically decreasing.
The maturity $\mu(\coord, t) = \iota(\tau(\coord, t))$ can only
decrease (by Lemma~\ref{lem:order-pres}), so
$\DebtScore(\coord, t)$ is non-decreasing.  Since decay is
exponential, the sum of all debt scores is bounded below by the
expression above.  The limit follows from $e^{-\lambda_\coord t}
\to 0$.
\end{proof}

\subsection{Equilibrium Conditions}\label{subsec:equilibrium}

\begin{definition}[Debt Equilibrium]\label{def:equilibrium}
A \emph{debt equilibrium} is a state where $\dot{D}(t) = 0$, i.e.,
the rate of trust decay equals the rate of repair:
\[
  \Lambda \, \mathbf{T}(t) = \mathbf{R}(t).
\]
\end{definition}

\begin{corollary}[Minimum Repair Rate]\label{cor:min-repair}
To maintain debt at or below threshold $\DebtThreshold$, the repair
rate must satisfy $\|\mathbf{R}(t)\|_1 \geq
\|\Lambda \, \mathbf{T}_{\min}\|_1$, where $\mathbf{T}_{\min}$ is
the trust vector corresponding to threshold maturity.
\end{corollary}

\subsection{Worked Example: Temporal Decay with TrustAlgebra}%
\label{subsec:temporal-example}

\begin{lstlisting}[style=jugeo-python]
from jugeo.evidence.trust import TrustAlgebra, TrustLevel
import math

ta = TrustAlgebra()

# Simulate temporal decay over 12 months
initial_trust = TrustLevel.SOLVER_DISCHARGED
decay_rate = 0.15  # lambda per month

for month in range(13):
    # Compute effective trust via attenuation
    decay_factor = math.exp(-decay_rate * month)
    if decay_factor > 0.8:
        current = TrustLevel.SOLVER_DISCHARGED
    elif decay_factor > 0.5:
        current = TrustLevel.RUNTIME_WITNESSED
    elif decay_factor > 0.2:
        current = TrustLevel.COPILOT_SUGGESTED
    else:
        current = TrustLevel.ORACLE_PROPOSED

    attenuated = ta.attenuate(current, initial_trust)
    is_ok = ta.is_admissible(attenuated,
                             TrustLevel.SOLVER_DISCHARGED)
    print(f"Month {month:2d}: trust={attenuated}, "
          f"admissible={is_ok}")
\end{lstlisting}

\begin{lstlisting}[style=jugeo-python]
# Benchmark trust to calibrate decay rates
from jugeo.evidence.trust import TrustAlgebra, TrustLevel

ta = TrustAlgebra()

# Compare pre-decay and post-decay trust
pre = TrustLevel.SOLVER_DISCHARGED
post = TrustLevel.COPILOT_SUGGESTED

benchmark = ta.benchmark_trust(post, pre)
print(f"Trust gap: {benchmark}")
print(f"Degraded: {ta.compare(post, pre)}")

# Check if bottom is reached
print(f"At bottom: {post == ta.bottom()}")
print(f"Top trust: {ta.top()}")
\end{lstlisting}

\begin{definition}[Monotone Fixpoint of Zero-Debt States]%
\label{def:fixpoint}
Let $F : 2^{\Ob(\Site)} \to 2^{\Ob(\Site)}$ map a set of zero-debt
coordinates $Z$ to
\[
  F(Z) = Z \cup \{\coord \mid \mathrm{deps}(\coord) \subseteq Z
  \text{ and } \mu(\coord) \geq \mu^*(\coord)\}.
\]
Then $F$ is monotone on the power set lattice, and by the
Knaster--Tarski theorem, $F$ has a least fixpoint
$\mathrm{lfp}(F) = \bigcap \{Z \mid F(Z) \subseteq Z\}$.
This fixpoint characterizes the maximal set of coordinates that can
be rendered debt-free without external verification.
\end{definition}

\begin{theorem}[Fixpoint Characterization]\label{thm:fixpoint}
A site is globally debt-free if and only if
$\mathrm{lfp}(F) = \Ob(\Site)$.
\end{theorem}

\begin{proof}
The forward direction is immediate: if every coordinate has zero
debt, then $\Ob(\Site) \subseteq F(\Ob(\Site))$, so the fixpoint
is $\Ob(\Site)$.  Conversely, if $\mathrm{lfp}(F) = \Ob(\Site)$,
then every coordinate either has zero debt directly or has all
dependencies in the zero-debt set, which by induction along the
site's acyclic dependency structure implies $\DebtScore(\coord)=0$
for all $\coord$.
\end{proof}


% =====================================================================
\section{Maturity Assessment}\label{sec:maturity}
% =====================================================================

\begin{algorithm}[H]
\caption{Maturity Assessment}
\label{alg:maturity}
\begin{algorithmic}[1]
\STATE \textbf{Input:} Site $\Site$, target maturity $\mu^*$
\STATE \textbf{Output:} Debt report $R$
\STATE $R \gets \emptyset$
\FOR{each $\coord \in \Ob(\Site)$}
  \STATE $\tau \gets \text{current\_trust}(\coord)$
  \STATE $\mu(\coord) \gets \text{trust\_to\_maturity}(\tau)$
  \STATE $d \gets \max(0, \mu^*(\coord) - \mu(\coord))$
  \STATE $\text{eff} \gets \min(\mu(\coord),
         \min_{\coord'} \mu(\coord'))$ over dependencies
  \STATE $R \gets R \cup \{(\coord, \mu(\coord), \text{eff}, d)\}$
\ENDFOR
\RETURN $R$ sorted by debt score descending
\end{algorithmic}
\end{algorithm}

\begin{lstlisting}[style=jugeo-python]
from jugeo.evidence.trust import TrustLevel

MATURITY = {
    TrustLevel.CONTRADICTED: 0,
    TrustLevel.UNVERIFIED: 1,
    TrustLevel.COPILOT_SUGGESTED: 2,
    TrustLevel.ORACLE_PROPOSED: 3,
    TrustLevel.HUMAN_ATTESTED: 3,
    TrustLevel.RUNTIME_WITNESSED: 4,
    TrustLevel.SOLVER_DISCHARGED: 5,
    TrustLevel.MECHANICALLY_VERIFIED: 5,
}

def debt_score(current: TrustLevel, target: int) -> int:
    return max(0, target - MATURITY[current])

print(debt_score(TrustLevel.COPILOT_SUGGESTED, target=4))  # 2
print(debt_score(TrustLevel.SOLVER_DISCHARGED, target=4))  # 0
\end{lstlisting}

% =====================================================================
\section{Sheaf-Cohomological Debt}\label{sec:cohomological-debt}
% =====================================================================

We connect technical debt to sheaf cohomology, showing that the
first cohomology group $\Hcoh^1(\Site, \mathcal{D})$ classifies
irreducible debt clusters---groups of coordinates whose debt cannot
be repaired independently.

\subsection{The Debt Sheaf}\label{subsec:debt-sheaf}

\begin{definition}[Debt Sheaf]\label{def:debt-sheaf}
The \emph{debt sheaf} $\mathcal{D}$ on $\Site$ assigns to each open
$U \subseteq \Ob(\Site)$ the abelian group
\[
  \mathcal{D}(U) = \bigoplus_{\coord \in U} \mathbb{Z} \cdot
  \DebtScore(\coord),
\]
with restriction maps induced by the site structure:
$\res^U_V : \mathcal{D}(U) \to \mathcal{D}(V)$ is the canonical
projection for $V \subseteq U$.
\end{definition}

\begin{definition}[\v{C}ech Cohomology of Debt]\label{def:cech-debt}
Let $\mathfrak{U} = \{U_i\}_{i \in I}$ be an open cover of $\Site$.
Define the \v{C}ech complex:
\begin{align*}
  C^0(\mathfrak{U}, \mathcal{D}) &= \prod_{i} \mathcal{D}(U_i), \\
  C^1(\mathfrak{U}, \mathcal{D}) &= \prod_{i < j}
    \mathcal{D}(U_i \cap U_j),
\end{align*}
with coboundary $\delta^0 : C^0 \to C^1$ given by
$(\delta^0 s)_{ij} = \res^{U_j}_{U_i \cap U_j}(s_j) -
\res^{U_i}_{U_i \cap U_j}(s_i)$.
The first cohomology is $\Hcoh^1(\Site, \mathcal{D}) =
\ker(\delta^1) / \mathrm{im}(\delta^0)$.
\end{definition}

\begin{theorem}[Cohomological Obstruction]\label{thm:cohom-obstruct}
A class $[\alpha] \in \Hcoh^1(\Site, \mathcal{D})$ is nonzero if
and only if there exists a cluster of coordinates whose debt scores
cannot be independently reduced to zero---they form an
\emph{irreducible debt cluster}.
\end{theorem}

\begin{proof}[Proof sketch]
A nonzero class $[\alpha] \in \Hcoh^1$ means the local debt
reductions $s_i \in \mathcal{D}(U_i)$ fail to glue to a global
section on $\bigcup U_i$.  Concretely, the restriction maps encode
dependency constraints: repairing $\coord_a \in U_i \cap U_j$
requires consistent trust levels from both $U_i$ and $U_j$.  When
$[\alpha] \neq 0$, no such consistent global repair exists without
first resolving the underlying structural coupling.  The converse
follows from the fact that $\Hcoh^1 = 0$ implies every cocycle is a
coboundary, so local repairs glue globally.
\end{proof}

\begin{corollary}[Global Repairability]\label{cor:global-repair}
If $\Hcoh^1(\Site, \mathcal{D}) = 0$, then the debt is globally
repairable: every coordinate's debt can be reduced to zero by a
sequence of independent local repair actions.
\end{corollary}

\begin{proof}
$\Hcoh^1 = 0$ means every local repair cocycle extends to a global
section, i.e., the sheaf condition $\glue$ holds for debt repair
strategies.  By the descent condition $\descent$ of $\Site$, the
local repairs assemble into a global repair plan.
\end{proof}

\begin{remark}[Computing $\Hcoh^1$]\label{rem:compute-cohom}
For a site with $n$ coordinates and $m$ dependency edges, computing
$\Hcoh^1(\Site, \mathcal{D})$ reduces to linear algebra over
$\mathbb{Z}$: the coboundary maps are integer matrices of dimension
at most $m \times n$, and the quotient can be computed via Smith
normal form in $O(n^2 m)$ time.  In practice, most codebases yield
$\Hcoh^1 = 0$ (globally repairable) unless circular dependencies
are present.
\end{remark}

\subsection{Detecting Irreducible Clusters}%
\label{subsec:irreducible-clusters}

\begin{lstlisting}[style=jugeo-python]
from jugeo.geometry.site import SiteBuilder, Coordinate
from jugeo.geometry.site import CoordinateKind, MorphismKind
from jugeo.evidence.trust import TrustAlgebra, TrustLevel

def detect_irreducible_clusters(site):
    """Detect debt clusters via cohomological analysis."""
    ta = TrustAlgebra()
    # Build adjacency from morphisms
    adj = {}
    for m in site.morphisms:
        adj.setdefault(m.source.name, []).append(
            m.target.name)
        adj.setdefault(m.target.name, []).append(
            m.source.name)

    # Find connected components with positive debt
    visited = set()
    clusters = []
    for coord in site.coordinates:
        if coord.name in visited:
            continue
        if ta.compare(coord.trust_tier,
                      TrustLevel.SOLVER_DISCHARGED):
            visited.add(coord.name)
            continue  # no debt
        # BFS to find cluster
        cluster = set()
        queue = [coord.name]
        while queue:
            c = queue.pop(0)
            if c in cluster:
                continue
            cluster.add(c)
            for neighbor in adj.get(c, []):
                if neighbor not in visited:
                    queue.append(neighbor)
        visited.update(cluster)
        if len(cluster) > 1:
            clusters.append(cluster)

    # Check sheaf condition for each cluster
    irreducible = []
    for cluster in clusters:
        if not ta.sheaf_condition_check(
                [c for c in site.coordinates
                 if c.name in cluster]):
            irreducible.append(cluster)
    return irreducible
\end{lstlisting}


% =====================================================================
\section{The Debt Presheaf as a Natural Transformation}%
\label{sec:debt-nat-trans}
% =====================================================================

We now formalize the relationship between current trust and target
trust as a natural transformation between presheaves, providing a
category-theoretic characterization of technical debt.

\subsection{Presheaf Setup}\label{subsec:presheaf-setup}

\begin{definition}[Current Trust Presheaf]\label{def:current-presheaf}
The \emph{current trust presheaf}
$\mathcal{T}_{\mathrm{cur}} : \Site^{\mathrm{op}} \to \Talg$ assigns
to each coordinate $\coord$ its current trust level
$\mathcal{T}_{\mathrm{cur}}(\coord) = \tau(\coord)$, with restriction
maps
\[
  \res^{\coord_1}_{\coord_2} : \mathcal{T}_{\mathrm{cur}}(\coord_1)
  \to \mathcal{T}_{\mathrm{cur}}(\coord_2)
\]
defined by $\res^{\coord_1}_{\coord_2}(\tau) = \tau \tjoin
\tau(\coord_2)$ for each morphism $\coord_2 \to \coord_1$.
\end{definition}

\begin{definition}[Target Trust Presheaf]\label{def:target-presheaf}
The \emph{target trust presheaf}
$\mathcal{T}^* : \Site^{\mathrm{op}} \to \Talg$ assigns to each
coordinate $\coord$ the minimum trust required by its target
maturity: $\mathcal{T}^*(\coord) = \iota^{-1}(\mu^*(\coord))$,
where $\iota^{-1}$ takes the minimum trust level mapping to the
given maturity.
\end{definition}

\subsection{The Debt Natural Transformation}%
\label{subsec:debt-nat-trans}

\begin{definition}[Debt Transformation]\label{def:debt-transform}
The \emph{debt transformation} is a family of maps
\[
  \eta_\coord : \mathcal{T}_{\mathrm{cur}}(\coord) \to
  \mathcal{T}^*(\coord), \qquad
  \eta_\coord(\tau) = \tau \tjoin \mathcal{T}^*(\coord),
\]
indexed by $\coord \in \Ob(\Site)$.
\end{definition}

\begin{theorem}[Naturality]\label{thm:naturality}
The family $\eta = (\eta_\coord)_{\coord \in \Ob(\Site)}$ is a
natural transformation
$\eta : \mathcal{T}_{\mathrm{cur}} \Rightarrow \mathcal{T}^*$.
That is, for every morphism $f : \coord_2 \to \coord_1$ in $\Site$,
the following diagram commutes:
\[
\begin{tikzcd}
  \mathcal{T}_{\mathrm{cur}}(\coord_1)
    \arrow[r, "\eta_{\coord_1}"]
    \arrow[d, "\res^{\coord_1}_{\coord_2}"']
  & \mathcal{T}^*(\coord_1)
    \arrow[d, "\res^{*\,\coord_1}_{\coord_2}"] \\
  \mathcal{T}_{\mathrm{cur}}(\coord_2)
    \arrow[r, "\eta_{\coord_2}"]
  & \mathcal{T}^*(\coord_2)
\end{tikzcd}
\]
\end{theorem}

\begin{proof}
We must show $\res^{*\,\coord_1}_{\coord_2} \circ \eta_{\coord_1}
= \eta_{\coord_2} \circ \res^{\coord_1}_{\coord_2}$.

Starting from $\tau \in \mathcal{T}_{\mathrm{cur}}(\coord_1)$:
\begin{align*}
  (\res^{*\,\coord_1}_{\coord_2} \circ \eta_{\coord_1})(\tau)
  &= \res^{*\,\coord_1}_{\coord_2}
     \bigl(\tau \tjoin \mathcal{T}^*(\coord_1)\bigr) \\
  &= (\tau \tjoin \mathcal{T}^*(\coord_1)) \tjoin
     \mathcal{T}^*(\coord_2).
\end{align*}
From the other path:
\begin{align*}
  (\eta_{\coord_2} \circ \res^{\coord_1}_{\coord_2})(\tau)
  &= \eta_{\coord_2}(\tau \tjoin \tau(\coord_2)) \\
  &= (\tau \tjoin \tau(\coord_2)) \tjoin \mathcal{T}^*(\coord_2).
\end{align*}
These are equal by associativity and commutativity of $\tjoin$,
together with the monotonicity assumption
$\mathcal{T}^*(\coord_1) \tleq \mathcal{T}^*(\coord_2)$
when $\coord_2 \to \coord_1$ (the target respects the site structure).
\end{proof}

\begin{corollary}[Restriction Compatibility]\label{cor:restrict-compat}
The component maps $\eta_\coord$ are restriction-compatible:
for any covering family $\{f_i : \coord_i \to \coord\}$,
the debt at $\coord$ is determined by the debts at the covering
coordinates via the gluing axiom.
\end{corollary}

\begin{definition}[Debt Kernel and Image]\label{def:debt-kernel}
The \emph{debt kernel} at $\coord$ is
$\ker(\eta_\coord) = \{\tau \mid \eta_\coord(\tau) =
\mathcal{T}^*(\coord)\}$---the set of trust levels already meeting
the target.  The \emph{debt image} is
$\mathrm{im}(\eta_\coord) = \{\eta_\coord(\tau) \mid \tau \in
\mathcal{T}_{\mathrm{cur}}(\coord)\}$---the achievable target trust
levels.  A coordinate has zero debt iff
$\tau(\coord) \in \ker(\eta_\coord)$.
\end{definition}



% =====================================================================
\section{Debt-Aware Repair}\label{sec:repair}
% =====================================================================

\begin{theorem}[Repair Convergence]\label{thm:repair-convergence}
Let $\rho = (a_1, \ldots, a_n)$ be a repair plan where each action
either succeeds (promoting trust) or fails (leaving trust unchanged).
Then $D(\Site)$ is monotonically non-increasing under $\rho$:
$D(\Site) \geq D(\Site\text{ after }a_1) \geq \cdots \geq
D(\Site\text{ after }a_n)$.
\end{theorem}

\begin{proof}
Each successful $a_i$ at $\coord_i$ promotes $\mu(\coord_i)$ by at
least one level, reducing $\DebtScore(\coord_i)$.  Failed actions
leave $\mu$ unchanged.  Theorem~\ref{thm:degrade-prop} ensures
promotions do not degrade dependents.
\end{proof}

Coordinates are prioritized by $\mathrm{priority}(\coord) =
\alpha \cdot \DebtScore(\coord) + \beta \cdot w(\coord) +
\gamma \cdot |\mathrm{deps}(\coord)|$.

\begin{lstlisting}[style=jugeo-python]
def plan_repairs(reports, alpha=1.0, beta=0.5, gamma=0.3):
    """Generate prioritized repair plan."""
    scored = []
    for r in reports:
        if r.debt_score == 0:
            continue
        p = (alpha * r.debt_score + beta * r.criticality
             + gamma * len(r.dependents))
        scored.append((p, r))
    scored.sort(reverse=True)
    return [(r.coord, choose_method(r)) for _, r in scored]
\end{lstlisting}

% =====================================================================
\section{Repair Strategy Algorithms}\label{sec:repair-algorithms}
% =====================================================================

Building on the convergence guarantee of
Theorem~\ref{thm:repair-convergence}, we present two concrete repair
algorithms: a greedy strategy and an optimal batch strategy.

\subsection{Greedy Repair with Priority Queue}%
\label{subsec:greedy-repair}

\begin{algorithm}[H]
\caption{Greedy Debt Repair}
\label{alg:greedy-repair}
\begin{algorithmic}[1]
\STATE \textbf{Input:} Site $\Site$, target $\mu^*$, weights $w$,
  parameters $\alpha, \beta, \gamma$
\STATE \textbf{Output:} Repaired site $\Site'$ with $D(\Site') = 0$
  or minimal residual
\STATE Initialize priority queue $Q \gets \emptyset$
\FOR{each $\coord \in \Ob(\Site)$ with $\DebtScore(\coord) > 0$}
  \STATE $p \gets \alpha \cdot \DebtScore(\coord) + \beta \cdot
    w(\coord) + \gamma \cdot |\mathrm{deps}(\coord)|$
  \STATE Insert $(\coord, p)$ into $Q$
\ENDFOR
\WHILE{$Q \neq \emptyset$}
  \STATE $(\coord, p) \gets \text{ExtractMax}(Q)$
  \STATE $\text{method} \gets \text{SelectRepairMethod}(\coord)$
  \STATE $\text{result} \gets \text{AttemptRepair}(\coord,
    \text{method})$
  \IF{result = \textsc{success}}
    \STATE Update $\mu(\coord) \gets \mu(\coord) + 1$
    \IF{$\DebtScore(\coord) > 0$}
      \STATE Re-insert $\coord$ with updated priority
    \ENDIF
    \FOR{each dependent $\coord'$ of $\coord$}
      \STATE Update effective maturity of $\coord'$
      \STATE Recompute priority of $\coord'$ in $Q$
    \ENDFOR
  \ENDIF
\ENDWHILE
\RETURN $\Site$
\end{algorithmic}
\end{algorithm}

\begin{theorem}[Greedy Convergence]\label{thm:greedy-conv}
Algorithm~\ref{alg:greedy-repair} terminates in at most
$\|\nabla D\|_1$ iterations and produces a site with
$D(\Site') \leq D(\Site) - k$, where $k$ is the number of
successful repairs.
\end{theorem}

\begin{proof}
Each successful repair promotes exactly one coordinate by one
maturity level, reducing aggregate debt by at least $w(\coord) \geq 1$.
Since $D(\Site) < \infty$ and each iteration either succeeds
(reducing debt) or fails (removing $\coord$ from $Q$ when no
further repair is possible), the algorithm terminates.
By Lemma~\ref{lem:repair-lower}, at most $\|\nabla D\|_1$
successful repairs suffice.
\end{proof}

\subsection{Optimal Batch Repair via Dynamic Programming}%
\label{subsec:dp-repair}

When repair actions have costs, we seek the minimum-cost repair plan
achieving zero debt.

\begin{algorithm}[H]
\caption{Optimal Batch Repair (Dynamic Programming)}
\label{alg:dp-repair}
\begin{algorithmic}[1]
\STATE \textbf{Input:} Coordinates $\coord_1, \ldots, \coord_n$,
  debt vector $\nabla D$, cost function
  $\text{cost}(\coord, \ell_1, \ell_2)$ for promoting from level
  $\ell_1$ to $\ell_2$, budget $B$
\STATE \textbf{Output:} Repair plan $\rho^*$ maximizing debt
  reduction within budget
\STATE $\textit{dp}[0][0] \gets 0$ \COMMENT{Base case}
\FOR{$i = 1$ to $n$}
  \FOR{$b = 0$ to $B$}
    \STATE $\textit{dp}[i][b] \gets \textit{dp}[i-1][b]$
    \COMMENT{Skip $\coord_i$}
    \FOR{$\ell = \mu(\coord_i)+1$ to $\mu^*(\coord_i)$}
      \STATE $c \gets \text{cost}(\coord_i, \mu(\coord_i), \ell)$
      \STATE $r \gets w(\coord_i) \cdot (\ell - \mu(\coord_i))$
      \IF{$b \geq c$ \AND $\textit{dp}[i-1][b-c] + r >
        \textit{dp}[i][b]$}
        \STATE $\textit{dp}[i][b] \gets
          \textit{dp}[i-1][b-c] + r$
        \STATE Record: promote $\coord_i$ to level $\ell$
      \ENDIF
    \ENDFOR
  \ENDFOR
\ENDFOR
\STATE Backtrack through $\textit{dp}$ to extract $\rho^*$
\RETURN $\rho^*$
\end{algorithmic}
\end{algorithm}

\begin{theorem}[DP Optimality]\label{thm:dp-optimal}
Algorithm~\ref{alg:dp-repair} computes the maximum debt reduction
achievable within budget $B$ in $O(n \cdot B \cdot 5)$ time, where
$5$ is the number of maturity levels.
\end{theorem}

\begin{proof}
This is a variant of the bounded knapsack problem.  Each coordinate
$\coord_i$ has at most $4$ promotion options (levels $2$--$5$), each
with a cost and a debt reduction value.  The dynamic programming
recurrence correctly enumerates all feasible combinations, and the
optimality follows from the principle of optimal substructure.
\end{proof}

\subsection{Implementation with CyclicSystemCoordinator}%
\label{subsec:cyclic-repair}

\begin{lstlisting}[style=jugeo-python]
from jugeo.maturity.cyclic_picture\
    .the_system_should_be_cyclic_not_pi import (
    CyclicSystemCoordinator, CyclePhase,
    CycleRecord, CyclicSystemWitness, CycleMetrics,
)
from jugeo.orchestration.synthesis_orchestrator import (
    SynthesisOrchestrator,
)
from jugeo.evidence.trust import TrustAlgebra, TrustLevel

def execute_repair_plan(site, repair_plan):
    """Execute a repair plan using cyclic refinement."""
    ta = TrustAlgebra()
    coordinator = CyclicSystemCoordinator()
    orchestrator = SynthesisOrchestrator()

    results = []
    for coord_name, target_level in repair_plan:
        # Start a repair cycle for this coordinate
        record = CycleRecord(
            coordinate=coord_name,
            phase=CyclePhase.SPECIFY,
        )
        coordinator.begin_cycle(record)

        # Phase 1: Specify the target
        witness = CyclicSystemWitness(
            coordinate=coord_name,
            target_trust=target_level,
        )

        # Phase 2: Synthesize repair
        record.phase = CyclePhase.SYNTHESIZE
        synthesis = orchestrator.synthesize(
            coordinate=coord_name,
            witness=witness,
        )

        # Phase 3: Verify the repair
        record.phase = CyclePhase.VERIFY
        verified = ta.solver_verified(
            synthesis.result)
        new_trust = (TrustLevel.SOLVER_DISCHARGED
                     if verified
                     else TrustLevel.COPILOT_SUGGESTED)

        # Phase 4: Integrate
        record.phase = CyclePhase.INTEGRATE
        coordinator.complete_cycle(record)

        metrics = CycleMetrics(
            coordinate=coord_name,
            old_trust=ta.bottom(),
            new_trust=new_trust,
        )
        results.append(metrics)

    return results
\end{lstlisting}



% =====================================================================
\section{Lean 4 Proof Sketch}\label{sec:lean-proof}
% =====================================================================

Full proofs reside in \texttt{Paper68\_TechnicalDebt.lean}.

\begin{lstlisting}[style=jugeo-output]
-- Maturity levels as a finite lattice
inductive MaturityLevel where
  | initial | managed | defined | measured | optimizing
  deriving DecidableEq, Repr

instance : LE MaturityLevel where
  le a b := a.toNat <= b.toNat

-- Debt score is non-negative and bounded
theorem debt_score_bounded (mu mu_star : MaturityLevel) :
    0 <= debtScore mu mu_star /\
    debtScore mu mu_star <= 4 := by
  unfold debtScore; omega

-- Repair monotonically reduces debt
theorem repair_reduces_debt (s : Site) (c : Coord)
    (h : repairAction c = success) :
    aggregateDebt (applyRepair s c) <=
    aggregateDebt s := by
  unfold aggregateDebt applyRepair
  apply Finset.sum_le_sum
  intro c' _
  by_cases hc : c' = c
  · subst hc; simp [debtScore, h]; omega
  · simp [applyRepair, hc]

-- Degradation respects site structure
theorem degradation_presheaf_functorial
    (f : Mor s c1 c2) (d : Degradation c2) :
    restrict f (degradation c2 d) =
    degradation c1 (propagate f d) := by
  simp [restrict, degradation, propagate,
        trust_join_comm, trust_join_assoc]
\end{lstlisting}

% =====================================================================
\section{Implementation}\label{sec:implementation}
% =====================================================================

The technical debt subsystem comprises:
\begin{itemize}[noitemsep]
  \item \textbf{Maturity Assessor}: \MaturityAssessment\ in
        \texttt{jugeo/maturity.py}.
  \item \textbf{Trust Tracker}: \TrustPresheaf\ in
        \texttt{jugeo/trust\_presheaf.py}.
  \item \textbf{Degradation Monitor}: Watches dependency changes and
        triggers trust re-evaluation.
  \item \textbf{Repair Planner}: Prioritized repair plan generation.
\end{itemize}

\begin{lstlisting}[style=jugeo-python]
from jugeo import Site

site = Site.from_source("project/")
target = {"auth": 5, "utils": 3, "api": 4}
reports = site.maturity_assessment(target=target)

for r in reports:
    if r.debt_score > 0:
        print(f"{r.coord}: debt={r.debt_score}")

plan = site.plan_repairs(reports)
for coord, method in plan:
    result = site.orchestrate_verification(
        coord=coord, method=method)
    print(f"Repaired {coord}: {result.new_trust}")
\end{lstlisting}

% =====================================================================
\section{Practical Scope}\label{sec:case-studies}
% =====================================================================

The repository version audited for this paper does not expose stable
directory-wide maturity scoring, debt-gate CLIs, or scripted repair commands
of the form shown in earlier drafts. We therefore omit shell transcripts and
deployment claims here. At present, the trustworthy public surface is the core
trust algebra plus small scripted analyses over explicit benchmark outputs.



% =====================================================================
\section{Evaluation}\label{sec:evaluation}
% =====================================================================

\begin{table}[H]
\centering
\caption{Technical Debt Assessment Results}
\label{tab:debt-results}
\begin{tabular}{lrrr}
\toprule
\textbf{Metric} & \textbf{Value} & \textbf{Unit} \\
\midrule
Total Programs       & \ppLXVIIIprogramCount & programs \\
Clean Programs       & \ppLXVIIIcleanCount   & programs \\
Debt-laden Programs  & \ppLXVIIIdebtCount    & programs \\
Verified             & \ppLXVIIIverifiedPct  & \% \\
Total Propositions   & \ppLXVIIIpropsTotal   & props \\
Detected Bugs        & \ppLXVIIIbugsTotal    & bugs \\
Mean Runtime         & \ppLXVIIItimeMean     & per program \\
Total Time           & \ppLXVIIItimeTotal    & \\
\bottomrule
\end{tabular}
\end{table}

\begin{table}[H]
\centering
\caption{Observed Prototype Characteristics}
\label{tab:debt-by-size}
\begin{tabular}{lrrrr}
\toprule
\textbf{Subset} & \textbf{Count} & \textbf{Mean Coords} & \textbf{Mean Props}
  & \textbf{Mean Time} \\
\midrule
Clean       & \ppLXVIIIcleanCount & \ppLXVIIIcleanCoordsMean
  & \ppLXVIIIcleanPropsMean & \ppLXVIIIcleanTimeMean \\
Debt-laden  & \ppLXVIIIdebtCount  & \ppLXVIIIdebtCoordsMean
  & \ppLXVIIIdebtPropsMean  & \ppLXVIIIdebtTimeMean \\
\bottomrule
\end{tabular}
\end{table}

All programs in the Paper~68 corpus verified successfully, and no bugs were
reported. At the same time, the generated per-program trust and debt scores in
\texttt{data-paper68.tex} are constant (0.5 for every instance), so the
current prototype does \emph{not} empirically distinguish clean from
debt-laden examples.

\paragraph{Interpretation.}
The useful empirical signal in this dataset is runtime, not discrimination:
clean and debt-laden subsets have comparable site sizes and similar runtimes.
Any future debt-measurement claim should therefore be supported by regenerated
data that actually vary with program condition.

% =====================================================================
\section{Related Work}\label{sec:related}
% =====================================================================

\paragraph{Technical Debt Taxonomy.}
Cunningham~\cite{Cunningham1992} introduced the debt metaphor.
Kruchten et al.~\cite{Kruchten2012} and Ernst et al.~\cite{Ernst2015}
classified debt categories.  Martin Fowler's \emph{Technical Debt
Quadrant}~\cite{Fowler2009} distinguishes deliberate vs.\ inadvertent
and reckless vs.\ prudent debt.  Our framework captures all four
quadrants: deliberate debt appears as explicitly set low target
maturity, while inadvertent debt manifests as temporal decay.

\paragraph{Quantitative Debt Models.}
SQALE~\cite{Letouzey2012} and SonarQube~\cite{SonarQube} measure
debt via heuristic metrics based on code smells, duplications, and
complexity.  The CAST Software Analysis model~\cite{CAST2020}
extends this with structural quality metrics and automated sizing.
The Software Improvement Group (SIG) maintainability
model~\cite{Heitlager2007} benchmarks code against industry norms
using a star-rating system.  In contrast, our approach grounds debt
quantification in formal trust levels rather than syntactic heuristics,
enabling debt scores that reflect actual verification status.

\paragraph{SEI Technical Debt Taxonomy.}
The Software Engineering Institute's taxonomy~\cite{SEI2012}
classifies debt into architecture debt, build debt, code debt,
defect debt, design debt, documentation debt, infrastructure debt,
people debt, process debt, requirement debt, service debt, and test
automation debt.  Our maturity lattice subsumes several of these
categories: code debt and defect debt map to trust degradation,
while architecture debt corresponds to cohomological obstructions
($\Hcoh^1 \neq 0$).

\paragraph{Maturity Models.}
CMMI~\cite{CMMI2010} defines five process maturity levels.  We
adapt this structure to per-coordinate verification maturity within
sheaf theory.  The ISO/IEC 15504 (SPICE) process assessment
model provides a similar five-level scheme, which we extend with
formal trust annotations.

\paragraph{Proof Maintenance.}
Ringer et al.~\cite{Ringer2020} address evolving Coq proofs across
type equivalences.  Czajka and Kaliszyk~\cite{Czajka2018} study
machine learning for proof automation.  Our degradation presheaf
extends proof maintenance to a broader setting that includes
lightweight automated tracking of trust levels across arbitrary
code evolution, not just type changes.

\paragraph{Formal Methods and Debt.}
Leino~\cite{Leino2010} describes Dafny's approach to verified
software, where proof obligations serve a similar role to our trust
annotations.  The connection between formal methods and technical debt
was explored by Kiniry and Zimmerman~\cite{Kiniry2008}, who argue
that formal specifications can prevent debt accumulation.  Our
framework makes this connection precise through the debt presheaf
and its cohomological properties.

\paragraph{Economics of Technical Debt.}
Schmid~\cite{Schmid2013} models debt using real options theory.
Ampatzoglou et al.~\cite{Ampatzoglou2015} provide a systematic
literature review of financial approaches to debt management.  Our
approach complements these economic models by providing the
underlying formal metric (the debt score) that economic analyses
can be built upon.


% =====================================================================
\section{Conclusion}\label{sec:conclusion}
% =====================================================================

We presented a sheaf-theoretic framework for technical debt through
trust degradation and maturity assessment.  The debt score provides
a precise metric at the modeling level, but the current Paper~68 dataset only
supports a modest empirical conclusion: the prototype runs on
\ppLXVIIIprogramCount{} curated programs in \ppLXVIIItimeMean{} on average.
Because the generated debt scores are uniform and no bugs were reported, this
revision does not claim successful technical-debt discrimination. Future work
should regenerate the benchmark with non-degenerate scoring and then revisit
repair-planning claims empirically.

% ─────────────────────────────────────────────────────────────────────
\begin{thebibliography}{99}

\bibitem{Cunningham1992}
W.~Cunningham, ``The WyCash Portfolio Management System,''
OOPSLA Experience Report, 1992.

\bibitem{Kruchten2012}
P.~Kruchten, R.~Nord, and I.~Ozkaya, ``Technical Debt: From Metaphor
to Theory and Practice,'' IEEE Software, 2012.

\bibitem{Ernst2015}
N.~Ernst et al., ``Measure It? Manage It? Ignore It? Software
Practitioners and Technical Debt,'' FSE 2015.

\bibitem{Letouzey2012}
J.-L.~Letouzey, ``The SQALE Method for Evaluating Technical Debt,''
MTD Workshop, 2012.

\bibitem{CMMI2010}
CMMI Product Team, ``CMMI for Development, Version 1.3,''
Software Engineering Institute, 2010.

\bibitem{Ringer2020}
T.~Ringer et al., ``Proof Repair Across Type Equivalences,''
PLDI 2020.

\bibitem{SonarQube}
SonarSource, ``SonarQube: Continuous Code Quality,''
\url{https://www.sonarqube.org}, 2023.

\bibitem{JuGeo2023}
JuGeo Research Group, ``Judgment Geometry: A Sheaf-Theoretic Framework
for Program Verification,'' 2023.

\bibitem{Lean4}
L.~de~Moura et al., ``The Lean 4 Theorem Prover and Programming
Language,'' CADE 2021.

\bibitem{Z32008}
L.~de~Moura and N.~Bj{\o}rner, ``Z3: An Efficient SMT Solver,''
TACAS 2008.


\bibitem{Fowler2009}
M.~Fowler, ``Technical Debt Quadrant,''
\url{https://martinfowler.com/bliki/TechnicalDebtQuadrant.html}, 2009.

\bibitem{CAST2020}
CAST Software, ``Automated Software Analysis for Structural Quality,''
Technical Report, 2020.

\bibitem{Heitlager2007}
I.~Heitlager, T.~Kuipers, and J.~Visser, ``A Practical Model for
Measuring Maintainability,'' QUATIC 2007.

\bibitem{SEI2012}
Software Engineering Institute, ``A Field Study of Technical Debt,''
CMU/SEI-2012-SR-013, 2012.

\bibitem{Czajka2018}
L.~Czajka and C.~Kaliszyk, ``Hammer for Coq: Automation for
Dependent Type Theory,'' JAR 2018.

\bibitem{Leino2010}
K.~R.~M.~Leino, ``Dafny: An Automatic Program Verifier for
Functional Correctness,'' LPAR 2010.

\bibitem{Kiniry2008}
J.~R.~Kiniry and D.~R.~Zimmerman, ``Secret Ninja Formal Methods,''
FM 2008.

\bibitem{Schmid2013}
K.~Schmid, ``A Formal Approach to Technical Debt Decision Making,''
QoSA 2013.

\bibitem{Ampatzoglou2015}
A.~Ampatzoglou et al., ``The Financial Aspect of Managing Technical
Debt: A Systematic Literature Review,'' IST 2015.

\end{thebibliography}

\end{document}
